% !TEX program = xelatex
%
% KDF: Knowledge Decay Framework — arXiv preprint (v0.3, 2026-04-19)
%
% Build pipeline (arxiv / TeX Live 2026 compatible):
%   1. Regenerate paper_body.tex from paper.md (body only, from "## 1. Introduction"):
%        sed -n '/^## 1\. Introduction/,$p' paper.md > _body.md
%        pandoc _body.md -o paper_body.tex --shift-heading-level-by=-1 --wrap=preserve
%        rm _body.md
%   2. Compile:
%        xelatex paper.tex && bibtex paper && xelatex paper.tex && xelatex paper.tex
%
% arxiv's AutoTeX uses TeX Live and respects the %!TEX program magic comment,
% so uploading paper.tex + paper_body.tex + references.bib works directly.

\documentclass[11pt,a4paper]{article}

% ---- Unicode + fonts (XeLaTeX) ----
% Tectonic / TeX Live 2024+ supports fontspec; Latin Modern Roman ships with
% every standard TeX Live install, so this is arxiv-safe.
\usepackage{fontspec}
\setmainfont{Latin Modern Roman}
% For pdflatex fallback (limited Unicode coverage — our paper.md uses θ, ≤, ↔,
% so xelatex is required), the pdflatex-compatible alternative would be:
%   \usepackage[utf8]{inputenc}
%   \usepackage[T1]{fontenc}
%   \usepackage{lmodern}

% ---- Pandoc helper macros (auto-generated lists use \tightlist) ----
\providecommand{\tightlist}{%
  \setlength{\itemsep}{0pt}\setlength{\parskip}{0pt}}

% ---- Geometry + layout ----
\usepackage[margin=1in]{geometry}
\usepackage{microtype}
\usepackage{setspace}
\onehalfspacing

% ---- Math ----
\usepackage{amsmath,amssymb,amsthm}
\usepackage{mathtools}

% ---- Tables ----
\usepackage{booktabs}
\usepackage{array}
\usepackage{longtable}
\usepackage{tabularx}

% ---- Graphics + color ----
\usepackage{graphicx}
\usepackage{xcolor}

% ---- Hyperlinks ----
\usepackage[colorlinks=true,linkcolor=blue,urlcolor=blue,citecolor=blue]{hyperref}

% ---- Lists ----
\usepackage{enumitem}

% ---- Authors + title ----
\usepackage{authblk}

% ---- Theorem environments ----
\newtheorem{hypothesis}{Hypothesis}
\newtheorem{conjecture}{Conjecture}

% ---- Bibliography ----
\usepackage[numbers,square,sort&compress]{natbib}
\bibliographystyle{unsrtnat}

% ============================================================================
% Title + Authors
% ============================================================================

\title{KDF: A Deterministic Architecture for Finite-Resource
       Information Preservation\\
       \large Cross-Domain Evidence and Self-Refutation of Canonical Values}

\author{Yasuhiro Kuroki}
\affil{Independent researcher, Japan}
\affil{\texttt{garden.of.knowledge.chai@gmail.com}}
\affil{ORCID: \href{https://orcid.org/0009-0006-8943-9344}
                   {0009-0006-8943-9344}}

\date{v0.4 (v2.7 Phase 2 Retrospective integrated), 2026-05-10\\
      Patent JP 2026-027032 (filed 2026-02-24)\\
      Code: \href{https://github.com/ChaiCroquis/kdf-perovskite}
                 {github.com/ChaiCroquis/kdf-perovskite}
      (PolyForm Noncommercial 1.0.0; commercial license separate)\\
      Reproducibility: \href{https://github.com/ChaiCroquis/kdf-perovskite/blob/main/docs/REPRODUCIBILITY.md}
                            {docs/REPRODUCIBILITY.md}
      (F-001--F-100 reproducer index, Quick Start: F-072 / F-087 / F-097)}

% ============================================================================
% Document body
% ============================================================================

\begin{document}

\maketitle

% -----------------------------------------------------------------------------
% Abstract
% -----------------------------------------------------------------------------
\begin{abstract}
\noindent
\textbf{KDF (Knowledge Decay Framework)} is a deterministic graph-compression
technique: given a graph and a retention budget, selected nodes are preserved
\textbf{verbatim} while unselected nodes are discarded, distinguishing KDF from
content-transforming methods such as LLM fact extraction. KDF combines three
mechanisms---edge-based continuous-time exponential decay (metabolic control),
rarity protection under the absolute threshold $\deg_E(v) \le 1$, and
integrity discovery (analogy) via graph-Laplacian eigenvalue
fingerprints---a three-pillar structure that our related-work survey finds
independently recurring across ten disciplines, including mammalian memory
consolidation, immune clonal selection, Ginzburg-Landau critical phenomena,
continual-learning EWC, Hopfield associative memory, and K-SVD sparse
coding.

\medskip
Empirically, KDF delivers $7.7\times$ gain over industry-standard TTL on
LongMemEval (ICLR 2025), $F_1 = 0.747$ on a 2,182-note Obsidian vault
(Wilcoxon $p = 0.006$), $2.3\times$ over Random for rare-event preservation
on NASA HTTP logs, and a \textbf{$+3.06$-point rare-recall improvement in a
realistic streaming replay of the NASA log}---providing the first positive
empirical anchor for our narrowed thesis that dynamic control components find
their true use in streaming scenarios rather than static queries.

\medskip
However, we also transparently report refutations of our own prior claims.
The sandwich 2-threshold \textit{mechanism} is supported, but \textbf{our
patent's canonical values $(\theta_L, \theta_U) = (0.70, 0.80)$ are
empirically refuted across four benchmarks} (Hopfield mixture, direct
analogy, synthetic pairs, LoCoMo streaming): specific values require
domain-specific calibration. Additional honest negatives include OSS issue
generalization $\times 1.00$ across three repositories, paper rediscovery
$\times 0.83$, and Gaussian-Process inducing-point selection failures.

\medskip
\textbf{v2 update (2026-04-29 -- 2026-04-30, see Version 2 Addendum):}
Phase 2 / Phase 2.5 systematic testing produced a \textbf{7-pattern
empirical arc} (5 narrowing findings + 2 positive cross-domain
replications): the streaming benefit is narrowed to \textit{temporally
recurring rare events} (F-087 Apache replication: $-13.04$ pt sign
reversal); the canonical $\alpha = 2.0$ is narrowed to NASA-recurring-rare
specific (F-091); Claim 31 functional rare protection is narrowed to
non-adversarial settings (F-092); the \texttt{bias-detector} commercial
predictor path is withdrawn (F-090: $5/11 = 45.5\% \ll 70\%$ pre-registered
threshold); direct SOTA-beating positioning is retracted (F-073 / F-074 /
F-075 LOSS); and the recurring-rare benefit is empirically replicated
cross-domain on $N = 3$ datasets (NASA HTTP $+3.06$ pt + Apache error
$+3.67$ pt + BGL HW kernel $+33.33$ pt; F-094 / F-097 PASS) but
\textbf{FAILS} on distributed file system log (F-100 HDFS:
H\_R+\_literal $-23.08$ pt, H\_anomaly $-4.30$ pt) --- narrowing
cross-domain coverage to \textbf{web access log family + HW kernel log
family} specific. The v2 grounded position is four product candidates
(Obsidian / MovieLens niche / Mem0 hybrid Router / Git Core
preservation) plus the F-086 $\gamma$ domain-fit predictor.

\medskip
\noindent\textbf{Keywords:} information preservation, graph metabolism, rarity
protection, analogy discovery, Laplacian fingerprint, Ginzburg-Landau,
Hopfield attractor, Equitable Coreset Selection, complementary learning
systems, memory consolidation.
\end{abstract}

% -----------------------------------------------------------------------------
% Body — generated from paper.md via pandoc
% -----------------------------------------------------------------------------
%
% Run once to generate paper_body.tex:
%   pandoc paper.md -o paper_body.tex
%
% Then \input it here. The \input line below assumes paper_body.tex exists;
% if it does not, see BUILD.md for setup. The included file should contain the
% sections (§1 Introduction through Appendix B) in LaTeX form.

\section{Version 2 Addendum (2026-04-29) --- Phase 2 Retrospective}\label{version-2-addendum-2026-04-29-phase-2-retrospective}

This is \textbf{version 2} of the paper. Version 1 was deposited on Zenodo
(DOI: \href{https://doi.org/10.5281/zenodo.19651035}{10.5281/zenodo.19651035})
on 2026-04-19, anchored on F-072 NASA HTTP streaming. Phase 2 + Phase 2.5
(2026-04-20 -- 2026-04-30) added 23 systematic findings (F-073 -- F-100)
that \textbf{sharpen} rather than overturn v1 claims, producing a \textbf{7-pattern
empirical arc} (5 narrowing + 2 positive cross-domain replications). We
summarize the substantive changes here; the v1 body below is unmodified
except for short \texttt{{[}v2:\ ...{]}} -- \texttt{{[}v2.7:\ ...{]}} caveats inline at the
relevant claims, and a Phase 2 / Phase 2.5 subsection appended to §5
Empirical Evaluation. The comprehensive single-document retrospective
lives at \href{../PHASE_2_RETROSPECTIVE.md}{\texttt{docs/PHASE\_2\_RETROSPECTIVE.md}}.

\textbf{Substantive narrowing}:

\begin{itemize}
\tightlist
\item
  \textbf{F-087 (2026-04-29) --- streaming benefit narrows to temporally
  recurring rare.} Apache error log replication of the F-072 streaming
  protocol with a frequency-based rare definition (one-shot
  reconnaissance probes, freq \(\le 10\)) yielded \(-13.04\) pt vs.~C0
  static (sign reversal from F-072's \(+3.06\) pt). The Claim 14 streaming
  benefit is empirically valid only when rare events recur over time
  (e.g., NASA HTTP 4xx/5xx persistent failure modes). For one-shot rare
  events, decay washes out the signal and ActivationScore amplifies
  common-path dominance, inverting the desired ordering.
\item
  \textbf{F-091 (2026-04-29) --- Claim 10 (\(\alpha=2\) ``core of the invention'')
  narrows to NASA-recurring-rare specific.} \(\alpha_{\text{core}}\) sweep
  \(\{0.5, 1.0, 2.0, 3.0, 4.0\}\) on NASA + Apache streaming + MovieLens
  null control: NASA confirms \(\alpha=2.0\) is empirically optimal
  (diff \(0.00\) pt vs best, F-072 anchor reproduced bit-exact); Apache
  shows \(\alpha=4.0\) optimal, \(\alpha=2.0\) off by \(-17.39\) pt.~MovieLens
  null PASS (no decay path). Claim 10 mechanism supported, but the
  canonical \(\alpha=2.0\) is \textbf{not domain-universal}; one-shot rare
  requires domain-specific calibration.
\item
  \textbf{F-092 (2026-04-29) --- Claim 31 functional rare protection narrows
  to non-adversarial settings.} Real-data adversarial burst
  perturbation (NASA window 50 + 1000 rare-target events): controller
  mechanism PASSes boundedness (\(\alpha_{\text{edge}} \in [1.0, 2.5]\)
  all 100 windows) and recovery (diff \(0.0043\) at w55), but functional
  rare detection FAILs (recall\_perturbed \(= 0.000\) vs.~baseline
  \(0.4592\)). Adversarial degree inflation demotes natural rare resources
  from Rare to Core layer. Claim 31 controller-stability claim is
  empirically supported; the implicit rare-protection-under-attack
  reading is not. Production deploy requires upstream defense layer
  (rate-limit / provenance filter).
\end{itemize}

\textbf{Substantive withdrawals}:

\begin{itemize}
\tightlist
\item
  \textbf{Direct SOTA-beating positioning (F-073 / F-074 / F-075).} Phase 2
  Top 3 candidates (Wikipedia orphan / BGL anomaly / Citation
  interdisciplinary bridge) all produced LOSS results
  (\(-4.07\) / \(-12.92\) / complete loss respectively). KDF is empirically
  not a SOTA-beating direct competitor; it is a structural-niche
  complement, consistent with the ``decisive predictor'' framework
  already articulated in §6.1. Paper's narrow scope claim is
  \emph{strengthened} by these negatives.
\item
  \textbf{\texttt{bias-detector} byproduct as commercial applicability predictor
  (F-090).} Systematic test on N=21 datasets gave 5/11 = 45.5\% certain
  prediction accuracy, well below the 70\% pre-registered threshold. The
  v1 §6.1 / Abstract reference to ``Applicability is predictable a
  priori via a zero-dependency bias-detector metric'' is \textbf{withdrawn}.
  The F-086 γ domain-fit framework (separate from \texttt{bias-detector},
  using γ-check correlation rate) replaces it as the working predictor;
  see §5 Phase 2/2.5 subsection.
\end{itemize}

\textbf{Strengthening (v1 anchors confirmed under stress)}:

\begin{itemize}
\tightlist
\item
  \textbf{Mem0 + KDF Router} cross-model robustness reconfirmed
  (F-060: gpt-4o-mini + gpt-4.1-mini, \(+10\) -- \(+23\) pt date/time
  literal recall).
\item
  \textbf{MovieLens niche genre surfacing} validated across six genres
  (F-082 + F-085 Task 1: γ-check \(100\%\) Film-Noir / IMAX / Western /
  Musical / War / Documentary).
\item
  \textbf{Domain-fit predictor sharpened}: F-086 γ refines v1's ``decisive
  predictor'' with hub-peripheral / hub-biased / peer-network
  distinction, validated across N=5 meta-families (3 PASS / 2 REJECT).
\end{itemize}

\textbf{v2 grounded position}: Four product candidates with empirical anchors
(Obsidian / MovieLens niche / Mem0 hybrid / Git Core preservation) plus
the F-086 γ predictor. The v1 narrative arc (broad cross-domain hypothesis
→ decisive predictor → narrowed scope) is preserved; v2 adds a sharper
empirical boundary at the cost of two specific commercial positionings.

\textbf{v2 changelog (machine-extractable)}:

\begin{itemize}
\tightlist
\item
  2026-04-29: streaming benefit (Claim 14) narrowed to temporally
  recurring rare per F-087 Apache replication failure
\item
  2026-04-29: bias-detector commercial path withdrawn per F-090
  systematic test (45.5\% \textless{} 70\% threshold)
\item
  2026-04-29: Claim 10 (\(\alpha=2\) ``core of the invention'') narrowed
  to NASA-recurring-rare specific per F-091 cross-domain \(\alpha\) sweep
  (Apache optimal at \(\alpha=4.0\), \(\alpha=2.0\) off by \(-17.39\) pt)
\item
  2026-04-29: Claim 31 functional rare protection narrowed to
  non-adversarial settings per F-092 real-data burst perturbation
  (controller boundedness + recovery PASS, functional recall FAIL
  \(0.000 / 0.4592\))
\item
  2026-04-29: Phase 2 / Phase 2.5 subsection added to §5 with
  F-073 -- F-093 raw findings (12 → 15 rows)
\item
  2026-04-29 (v2.1): F-093 NASA F-072 anchor robustness to rare code
  subset --- anchor sharpening (not self-refutation): rare = 4xx/5xx is
  実質 404-pattern driven on this dataset (5xx response absent), V1/V2/V4
  identical \(+3.06\) pt confirms KDF captures rare \emph{resource pattern},
  not rare \emph{code set}. Generalization to 5xx-containing logs is future
  work
\item
  2026-04-29 (v2.2): F-094 Apache recurring-rare positive replication
  of F-072 streaming benefit --- H\_R+ PASS (\(\Delta_{\text{recurring}} =
  +3.67\) pt \(> +1.0\) pt threshold) on Apache same dataset with rare def
  flipped from one-shot (freq \(\le 10\), F-087 reference) to recurring
  (freq \(\ge 5\)). Sanity F-087 reproduce \(-13.04\) pt \(\pm 0.0\) pt
  (preprocessing / build env consistent). F-072 anchor's true axis
  (recurring rare structure, hypothesized via F-093 structural reading)
  is empirically replicated cross-domain (N=2: NASA \(+3.06\) pt + Apache
  recurring \(+3.67\) pt). 4-pattern self-refutation arc (F-070 / F-087 /
  F-091 / F-092, narrowing direction) is supplemented with \textbf{1 positive
  replication (F-094)} for arc completion (subsequently expanded to
  \textbf{7 patterns total = 5 narrow + 2 positive} by F-097 BGL HW kernel
  log positive replication and F-100 HDFS distributed file system log
  narrowing per v2.5 / v2.7 entries below). Scope of Claim 14 streaming
  benefit: recurring rare structure \$\times \alpha=2.0 \times \$
  web access log family + HW kernel log family
\item
  2026-04-29 (v2.3): Foreign baseline anchor (Mem0 family, N=3 empirical
  cells: F-048 local proxy / F-053 paid alone / F-060 paid hybrid Router)
  added to §5 --- setting-dependent position formalized, KDF as Mem0
  complementary layer
\item
  2026-04-29 (v2.4): F-095 / F-096 attempted local replication of F-060
  documented as honest infrastructure record. F-095 (\texttt{llama3.1:8b},
  default batching) infrastructure-infeasible at 5-Q checkpoint
  (\(\sim 33\) -- \(36\) h extrapolated wall-clock); F-096 (\texttt{qwen2.5:3b}
  with single-add / batch=50 ingest optimization, \(\sim 2.5\) h
  wall-clock, LongMemEval \(n=479\) + LoCoMo \(n=321\) both completed)
  inconclusive due to sub-noise floor (LoCoMo Mem0 alone \(0/321\), KDF
  alone \(2/321\), \(\Delta_{\text{router}} = +0.62\) pt mechanical PARTIAL,
  \(p = 0.5\)). Sanity 2 local-vs-paid baseline shift \(-20\) -- \(-64\) pt
  across cells indicates qwen2.5-3b is below the threshold for valid
  H\_R measurement on these benchmarks. F-095 / F-096 do \textbf{not refute
  F-060} --- they anchor the local-LLM infrastructure threshold; valid
  local replication requires \(7\)B+ local LLM, left as future work.
  Descriptive cross-LLM-size finding: Mem0 framework batched fact
  compression destroys \(\sim 99\%\) of raw turn substrings (recall
  \(0.000\) -- \(0.006\) on 8B 5-Q + 3B 479-Q both), sharpening F-048
  caveat from ``weak LLM degrades retrieval'' to ``framework compression
  itself does not preserve substrings independent of LLM size''
\item
  2026-04-29 (v2.5): F-097 BGL recurring-rare cross-domain N=3 expansion
  of F-094 --- H\_R+ PASS (\(\Delta_{\text{recurring}} = +33.33\) pt
  \(> +1.0\) pt threshold) on BGL Blue Gene/L HW kernel failure log
  (300,000-line subsample, anomaly subgraph 79,641 records, 8 unique
  content templates). recurring-rare benefit replicated cross-domain
  in N=3 settings (NASA HTTP + Apache error log + BGL HW kernel log),
  expanding the durable axis from web log family to HW kernel log
  family. \textbf{Sanity inconsistent} (V\_one-shot \(\Delta = +0.00\) pt,
  expected negative direction per F-087 / F-094 anchor): one-shot
  disposal mechanism narrowed to web log family with rich resource
  alphabet --- BGL alphabet is small (8 anomaly templates total) so
  top-30\% selection captures most one-shots regardless of streaming.
  4-pattern self-refutation arc + 1-positive arc (F-094 v2.2) is now
  \textbf{6-pattern arc} (4 narrow F-070/F-087/F-091/F-092 + 2 positive
  F-094 web + F-097 HW): cross-domain durability of recurring direction
  empirically anchored across web access logs and HW kernel logs.
  Caveat: \(n_{\text{rare}} = 3\) in BGL recurring set, recall granularity
  \(\{0, 1/3, 2/3, 1\}\) only; direction unambiguous, low-N power
\item
  2026-04-29 (v2.6): F-099 v1 router (precision-only, no length filter)
  characterization across 5 cells × 3 variants --- H\_v1\_paid PASS\_negative
  on both LongMemEval paid cells (F-053 \(\Delta_{v1}=-11.60\) pt
  \(p=10^{-7}\), F-059 \(\Delta_{v1}=-13.00\) pt \(p=10^{-10}\)); H\_v1\_local
  REPRODUCED (F-096 LongMemEval local \(\Delta_{v1}=+6.26\) pt deterministic
  post-hoc replication); Sanity 4/4 paid v2 cells match F-060 published
  values within \(\pm 0.0004\). \textbf{F-099 anchors v2's length\$\ge\$100
  component necessity} (indirect logic from v1 -11/-13 pt vs v2 0 pt
  on paid LongMemEval): without the length filter, routing precision
  queries to KDF (the lower performer in paid setting where Mem0 \(>>\)
  KDF) is actively harmful. Important caveat: \textbf{F-099 does NOT
  establish the precision component's independent necessity} --- the
  v3 variant (length-only, no precision filter) yields identical
  \(+9.66 / +22.43\) pt gain on LoCoMo cells as v2 (precision AND length),
  suggesting the precision filter may be redundant once the length
  filter is satisfied on these test cells. Precision filter necessity
  remains evaluated by other anchors (e.g., F-060 design rationale),
  not by F-099. v1 is \textbf{not a product candidate} (paid hurt significant,
  local help at sub-noise floor). Router design space now fully
  documented (5 cells × 3 variants). \textbf{First test case of pre-reg
  template \texttt{\_template\_pre\_reg.md}} (operational form of trigger 5 /
  post-hoc narrowing / observation-interpretation / recommendation-
  boundary protocols, introduced after Direction A occurrence 1-4
  showed body memory dependence as root failure mode). F-099 is a
  \textbf{low-severity} test case (post-hoc deterministic, no LLM, no GPU
  bench, no apples-to-apples confusion) --- §0 catches confirmed here
  are self-evident; \textbf{structural prevention of Direction A occurrence
  5 awaits a higher-severity F-100+ second test case} (LLM bench /
  wall-clock estimation / framework behavior comparison). Independent
  v1 verification on fresh paid data is also deferred to F-100+
\item
  2026-04-30 (v2.7): F-100 HDFS empirical verification --- cross-domain
  N=4 attempt FAILED on Loghub HDFS\_v1 distributed file system log
  (100,000-BlockId subsample, 2.35M edges, 26 templates, 4.90\% anomaly).
  H\_R+\_literal Δ=\(-23.08\) pt FAIL (recurring rare benefit narrowing
  refuted on HDFS), H\_anomaly Δ=\(-4.30\) pt FAIL (HDFS-native anomaly
  preservation FAIL), Sanity inconsistent (one-shot disposal absent,
  reinforcing F-097 small-alphabet caveat). Cross-domain durability
  arc narrowed to \textbf{web access log family + HW kernel log family}
  specific (NASA HTTP + Apache error log + BGL Blue Gene/L), distributed
  file system log NOT covered. 6-pattern arc → \textbf{7-pattern arc}
  (5 narrow + 2 positive). Pre-reg \texttt{\_template\_pre\_reg.md} second
  high-severity test case: §0.1 anchor constraint predicted small-
  alphabet caveat would re-confirm → empirically borne out (sanity
  +0 pt, identical to F-097 BGL pattern); §0.3 observation/interpretation
  separation prevented narrative drift; protocol §5 explicit禁止
  ``literal FAIL but H\_anomaly PASS as N=4'' not triggered (both FAIL).
  ONE\_PAGER.md Path A SIEM PoC scope adjusted: HDFS-style distributed
  file system log out of scope, web/HW kernel log domain remains valid
\item
  v1 body otherwise unmodified; inline \texttt{{[}v2:\ ...{]}} / \texttt{{[}v2.1:\ ...{]}} /
  \texttt{{[}v2.2:\ ...{]}} / \texttt{{[}v2.3:\ ...{]}} / \texttt{{[}v2.4:\ ...{]}} / \texttt{{[}v2.5:\ ...{]}} markers
  indicate caveat insertion points
\item
  \textbf{7-pattern arc anchors the paper's honesty-first epistemic stance}
  (5 narrowing: F-070 / F-087 / F-091 / F-092 / F-100; 2 positive
  cross-domain replication: F-094 web + F-097 HW): mechanism supported,
  specific application robustness narrowed, cross-domain durable on the
  recurring-rare axis across web access log family + HW kernel log
  family but NOT distributed file system log family
\item
  \textbf{F-093 adds anchor sharpening category} (distinct from
  self-refutation): empirical anchor's true scope identified at higher
  resolution by dataset structural investigation
\end{itemize}

\begin{center}\rule{0.5\linewidth}{0.5pt}\end{center}

\section{1. Introduction}\label{introduction}

\subsection{1.1 Problem}\label{problem}

In persistently growing information networks---LLM-agent conversation memory, personal knowledge bases, distributed-system logs, OSS issues, citation networks---storage constraints force a tradeoff between \textbf{volume} and \textbf{quality}:

\begin{itemize}
\tightlist
\item
  \textbf{Full retention} is practically infeasible at the TB/day scale.
\item
  \textbf{Random reduction} (random sampling, reservoir sampling) \textbf{statistically loses} rare-but-important items (e.g., 4xx/5xx error logs, forgotten-yet-consequential notes, minority-language utterances).
\item
  \textbf{Label-dependent methods} (stratified sampling, tail-based sampling, active learning, Equitable Coreset Selection) are powerful but presuppose labels or ground truth whose availability is not guaranteed in production.
\end{itemize}

\textbf{Requirement}: without labels, using only structure, protect rare items that will be needed later while metabolizing redundant information.

\begin{center}\rule{0.5\linewidth}{0.5pt}\end{center}

\subsection{1.2 Observation}\label{observation}

\textbf{Broader context.} The general-systems tradition, exemplified by ``Living Systems Theory'' \citep{miller1978livingsystems}, has long identified twenty ``critical subsystems''---\emph{associator}, \emph{memory}, \emph{decoder}, and others---shared across the seven hierarchical levels from cell to supranational system. Our observation is \textbf{consistent with} this tradition; we narrow our contribution to three of those twenty subsystems---the ones specialized for information preservation (metabolism, rarity protection, recombination)---and to presenting a concrete numerical implementation.

Having attempted to solve this problem through independent engineering effort, we then noticed that the three-pillar architecture we arrived at \textbf{corresponds structurally} to the following (a qualitative observation of parallelism, not a proof of mathematical isomorphism):

\begin{enumerate}
\def\labelenumi{\arabic{enumi}.}
\tightlist
\item
  \textbf{Mammalian memory consolidation} (Complementary Learning Systems) \citep{mcclelland1995cls}
\item
  \textbf{Memory B-cell selection in the immune system} (germinal-center affinity maturation)
\item
  \textbf{Ginzburg-Landau theory of critical phenomena} (quartic stabilization term)
\item
  \textbf{Elastic Weight Consolidation in continual learning} \citep{kirkpatrick2017ewc}
\item
  \textbf{Equitable Coreset Selection} \citep{wang2024ecs} and the unlearning variant \textbf{UPCORE} \citep{patil2025upcore}
\item
  \textbf{Hopfield associative memory networks} \citep{hopfield1982neural} and modern Hopfield networks \citep{ramsauer2021hopfield}
\item
  \textbf{Continuous-time Markov processes on graphs} (spectral-gap convergence)
\item
  \textbf{Tail-risk management for Pareto heavy-tailed distributions} \citep{pareto1896economie}
\item
  \textbf{Rare-atom preservation in K-SVD sparse dictionary learning} \citep{aharon2006ksvd}
\item
  \textbf{Two-factor synaptic consolidation} \citep{tartaglia2025twofactor}
\end{enumerate}

\begin{center}\rule{0.5\linewidth}{0.5pt}\end{center}

\subsection{1.3 Claims}\label{claims}

This paper advances three claims:

\begin{itemize}
\item
  \textbf{(C1) Integration}: the findings from the ten disciplines above admit a uniform description as a three-pillar structure---\textbf{metabolic control + rarity protection + recombination}. KDF is the first \textbf{numerical implementation} of this unified description.
\item
  \textbf{(C2) Structural similarity (not yet rigorous)}: two distinctive KDF formulas---(a) the meta-control law \(\Delta \alpha \propto \delta k^4\) and (b) the continuous-time kernel \(\lambda(C) \cdot \exp(-\lambda \cdot dt)\)---share \textbf{the same functional form} as, respectively, the quartic term of the Ginzburg-Landau free energy and the survival probability of a continuous-time Markov process. This is a qualitative structural correspondence; a proof as mathematical isomorphism (a structure-preserving bijection) remains future work.
\item
  \textbf{(C3) Novel contribution (mechanism supported; canonical values refuted across four benchmarks)}: the \textbf{sandwich 2-threshold mechanism}---acceptance only within a middle band, imposed jointly by a lower threshold \(\theta_L\) and an upper threshold \(\theta_U\)---has no counterpart in the ten surveyed disciplines and is a distinctive KDF contribution. However, the \textbf{specific values that patent claims 46/48 designate as canonical, \((\theta_L, \theta_U) = (0.70, 0.80)\), are empirically refuted across the following four benchmarks}:

  \begin{enumerate}
  \def\labelenumi{\arabic{enumi}.}
  \tightlist
  \item
    \textbf{Phase V3 Hopfield} (F-041): at \(\theta_U = 0.80\), \(0\%\) of Hopfield mixture states are detected; rejection of \(24\%\) (at \(P = 18\)) and up to \(\sim 40\%\) at higher load factors becomes feasible only after lowering the threshold to \(\theta = 0.40\).
  \item
    \textbf{F-068 analogy-discovery (direct)}: on Gentner classics and git↔paper cross-domain pairs, scores concentrate uniformly at \(\ge 0.99\), so an upper bound of \(0.80\) would reject every true positive.
  \item
    \textbf{F-070 Part A synthetic + F-068 scenarios}: over 38 pairs (22 positive / 16 negative), F1(canonical) = \(0.000\) whereas F1(\((0.70, 1.00)\)) = \(1.000\).
  \item
    \textbf{F-070 Part B LoCoMo Rev12 streaming}: across 30 queries and 1,140 RARE nodes, the canonical values cause 100\% demotion to Garbage after the 60-cycle timeout (complete information loss).
  \end{enumerate}

  We therefore narrow the claim to: ``\textbf{the mechanism is meaningful across domains; the specific values require domain-specific calibration.}'' §4.2 elaborates. We claim novelty for the \emph{mechanism} only, and we do not claim universality for the canonical numerical values.
\end{itemize}

\begin{center}\rule{0.5\linewidth}{0.5pt}\end{center}

\subsection{1.4 Tension between Universality and Novelty}\label{tension-between-universality-and-novelty}

The paper's assertion that ``KDF's structure has been independently found in ten disciplines'' is, from a novelty standpoint, \textbf{a double-edged sword}:

\begin{itemize}
\tightlist
\item
  \textbf{Favorable reading}: ten disciplines converging independently on the same solution suggests that KDF, as the first integrated implementation of this universal pattern, has genuine significance.
\item
  \textbf{Critical reading}: analogous solutions are already known in ten disciplines, so KDF is mere reinvention and lacks novelty.
\end{itemize}

We confront this tension directly and restrict our novelty claim to the following three narrow items:

\begin{enumerate}
\def\labelenumi{\arabic{enumi}.}
\tightlist
\item
  \textbf{Novelty of integration}: we find no prior instance of a single numerical implementation or a single patent claim that integrates all three pillars (metabolism, rarity protection, recombination). Each discipline specializes in one or two pillars (neuroscience focuses mainly on metabolism and recombination, ECS mainly on rarity protection, and so on).
\item
  \textbf{Novelty of the sandwich 2-threshold \emph{mechanism} (but not its specific values)}: the \emph{dual-threshold structure} that admits only a middle band via a lower bound \(\theta_L\) and an upper bound \(\theta_U\) has no counterpart in the ten disciplines we surveyed and is distinctive to KDF. However, the patent canonical values \((0.70, 0.80)\) have been empirically refuted across four benchmarks (§1.3 C3, §4.2); we therefore claim novelty only for the \emph{mechanism}, and we do not claim optimality for any specific numerical values.
\item
  \textbf{Novelty of an open numerical implementation}: this is the first publicly released implementation to provide concrete parameter values (\(\beta = 0.01\), 7:2:1, 5:3:1, \([0.70, 0.80]\), \(\delta k^4\)) together with a source-available implementation (PolyForm Noncommercial 1.0.0; commercial license available separately).
\end{enumerate}

We do \textbf{not} claim novelty for any individual pillar (synaptic pruning, EWC, Laplacian fingerprinting, and the like) in isolation. These are pre-existing contributions.

\begin{center}\rule{0.5\linewidth}{0.5pt}\end{center}

\subsection{1.5 Paper Organization}\label{paper-organization}

Section 2 presents KDF's three mechanisms and key formulas concisely. Section 3 systematizes the correspondences with the ten independent disciplines. Section 4 argues the three structural similarities. Section 5 presents empirical evidence (six positive and four negative cases). Section 6 discusses implications and limitations. Section 7 concludes with a refined applicability predictor and the 2 × 2 Mem0 / KDF benchmark matrix. Section 8 develops a stronger theoretical foundation by aligning KDF with Burt's Structural Holes theory, paired with an explicit Limitations and Risks subsection that makes the hierarchy vs.~§4 deliberate. Acknowledgments, References (generated from \texttt{references.bib}), and Appendices A (key formulas) and B (implementation architecture) follow.

\begin{center}\rule{0.5\linewidth}{0.5pt}\end{center}

\section{2. KDF Architecture Overview}\label{kdf-architecture-overview}

\subsection{2.1 Basic Structure}\label{basic-structure}

Let the information structure be a graph \(G = (V, E)\). Each edge \((u, v) \in E\) carries parameters: strength \(w_{uv}\), connection history \(n_{uv}\), and last-access time \(t^{\text{acc}}_{uv}\).

\textbf{Mechanism 1 (Metabolic Control).} Using the local congestion \(C_{uv} = \deg(u) + \deg(v)\), we define the decay rate

\[
\lambda(C_{uv}) = \beta \left( 1 + \gamma C_{uv}^{\alpha} \right)
\]

and apply exponential edge-weight decay over a discrete time step \(dt\):

\[
w_{uv}(t + dt) = w_{uv}(t) \cdot \exp\bigl(-\lambda(C_{uv}) \cdot dt\bigr) .
\]

Threshold-based pruning and probabilistic pruning with \(\Pr[\text{prune}] = 1 - \exp(-\lambda \cdot dt)\) are also permitted.

\textbf{Mechanism 2 (Rarity Protection).} Any node \(v\) satisfying the absolute threshold \(\deg_E(v) \le 1\) is treated as a rare object and \textbf{unconditionally excluded} from metabolic control over the protection interval

\[
[t_0, t_0 + T_{\text{wait1}}] \cup (t_0 + T_{\text{wait1}}, t_0 + T_{\text{wait1}} + T_{\text{wait2}}]
\]

(a two-stage review). The two intervals satisfy \(T_{\text{wait1}} = T_{\text{wait2}} \in [30, 70]\) --- Claim 37 requires equal lengths, Claim 39 specifies the range, and the canonical default is \(T_{\text{wait}} = 50\).

\textbf{Mechanism 3 (Integrity Discovery).} From the ego-subgraph of a rare node we construct the Laplacian \(L_v\), compute its eigenvalue vector \(\phi(v) \in \mathbb{R}^{32}\)---a fixed-length, isometry-invariant fingerprint---and evaluate the integrity score against existing nodes:

\[
S(v, u) = a \cdot S_{\text{cos}}(\phi(v), \phi(u)) + b \cdot S_{\text{struct}}(v, u) + c \cdot S_{\text{sign}}(v, u), \qquad a, b, c > 0 .
\]

The formula above defines the \textbf{inner} similarity \(S_{\text{inner}}\); the \textbf{outer} (aggregated) integrity score --- on which the sandwich acceptance condition actually operates --- is obtained by aggregating \(S_{\text{inner}}\)-derived values of three types (systematic, relational, attribute). A new edge is generated when the outer score satisfies the \textbf{sandwich 2-threshold criterion} \(\theta_L \le S_{\text{outer}} \le \theta_U\) (canonical defaults \(\theta_L = 0.70\), \(\theta_U = 0.80\)). The weights are used at two independent levels:

\begin{itemize}
\tightlist
\item
  \textbf{Inner fingerprint similarity} \(S_{\text{inner}}\) (combination of structural fingerprints): \(a : b : c = 0.40 : 0.35 : 0.25\), the linear combination of \(S_{\text{cos}}\), \(S_{\text{struct}}\), and \(S_{\text{sign}}\) (Claim 45).
\item
  \textbf{Outer integrity aggregation} \(S_{\text{outer}}\) (aggregation over systematic / relational / attribute similarities): \(\text{systematic} : \text{relational} : \text{attribute} = 7 : 2 : 1\) (Claim 44).
\end{itemize}

The two weightings are independent: the inner one governs similarity between fixed-length vectors, while the outer one governs the aggregation of classified similarity scores. The sandwich threshold is imposed on the outer score.

\begin{center}\rule{0.5\linewidth}{0.5pt}\end{center}

\subsection{2.2 Meta-control (Claims 27--32)}\label{meta-control-claims-2732}

Let \(\delta k = \max(0, \langle k \rangle - k_{\text{opt}})\) be the deviation of the network's mean degree \(\langle k \rangle\) from the target degree \(k_{\text{opt}}\). The adaptive law

\[
\Delta \alpha = -\eta (H - H_{\text{target}}) \pm \mu \cdot \delta k^{4}
\]

updates \(\alpha\) within the range \([\alpha_{\min}, \alpha_{\max}]\). The Lyapunov stability condition \(\eta^2 > \mu^2\) holds (full Lyapunov analysis is provided in the filed patent JP 2026-027032 and in the reference implementation \texttt{crates/cgb-kdf/src/meta\_control.rs}; not reproduced in this paper).

\begin{center}\rule{0.5\linewidth}{0.5pt}\end{center}

\subsection{2.3 Hierarchical Management Regions}\label{hierarchical-management-regions}

The architecture maintains three regions---short-term, long-term, and rare---with update-period ratio \(dt_1 : dt_2 : dt_3 = 5 : 3 : 1\).

\begin{center}\rule{0.5\linewidth}{0.5pt}\end{center}

\section{3. Three-Pillar Structure Across Ten Independent Disciplines}\label{three-pillar-structure-across-ten-independent-disciplines}

The table below enumerates the concrete realization of the three pillars in each discipline. The correspondence with KDF \textbf{can be observed at a qualitative level} (this is not a claim of strict equivalence).

{\def\LTcaptype{none} % do not increment counter
\begin{longtable}[]{@{}
  >{\raggedright\arraybackslash}p{(\linewidth - 8\tabcolsep) * \real{0.2000}}
  >{\raggedright\arraybackslash}p{(\linewidth - 8\tabcolsep) * \real{0.2000}}
  >{\raggedright\arraybackslash}p{(\linewidth - 8\tabcolsep) * \real{0.2000}}
  >{\raggedright\arraybackslash}p{(\linewidth - 8\tabcolsep) * \real{0.2000}}
  >{\raggedright\arraybackslash}p{(\linewidth - 8\tabcolsep) * \real{0.2000}}@{}}
\toprule\noalign{}
\begin{minipage}[b]{\linewidth}\raggedright
Discipline
\end{minipage} & \begin{minipage}[b]{\linewidth}\raggedright
Metabolism
\end{minipage} & \begin{minipage}[b]{\linewidth}\raggedright
Rarity protection
\end{minipage} & \begin{minipage}[b]{\linewidth}\raggedright
Recombination / integration
\end{minipage} & \begin{minipage}[b]{\linewidth}\raggedright
Representative references
\end{minipage} \\
\midrule\noalign{}
\endhead
\bottomrule\noalign{}
\endlastfoot
Mammalian brain (1) & Synaptic pruning during sleep & Engram plasticity window & Hippocampus → neocortex replay & \citet{mcclelland1995cls}; \citet{tartaglia2025twofactor} \\
Immune system (2) & Naïve B-cell removal & Memory B-cell pool (restricted clonality) & Germinal-center affinity maturation & \citet{mesin2020gc} \\
Critical-phenomena physics (3) & SOC avalanche decay & Power-law tail (critical cluster) & Ginzburg-Landau quartic stabilization & \citet{bak1987soc}; \citet{ginzburg1950superconductivity} \\
Continual-learning ML (4) & Weight decay / dropout & EWC Fisher protection & Replay buffer & \citet{kirkpatrick2017ewc} \\
Coreset selection (5) & Adaptive pruning & ECS minority preservation & Class-sensitive partitioning & \citet{sener2018coreset}; \citet{wang2024ecs} \\
Hopfield associative memory (6) & Dynamic pattern decay & Stored-pattern attractor basin & \textbf{Spurious-attractor rejection remains unsolved} & \citet{hopfield1982neural}; \citet{ramsauer2021hopfield} \\
Markov processes on graphs (7) & \(\exp(-\lambda \cdot dt)\) mixing & Slow-mixing mode (\(\lambda_2\)) & Spectral gap / Fiedler vector & \citet{levin2017markov} \\
Economics / finance (8) & Portfolio rebalance & Tail-event capital reserve & Pareto / extreme-value theory & \citet{pareto1896economie}; \citet{embrechts1997extreme} \\
Signal processing (9) & Dictionary pruning & Rare-atom preservation & K-SVD iterative update & \citet{aharon2006ksvd} \\
Graph ML (10) & Time-decayed line graph & Rare-node structural signal & Analogy discovery via fingerprint & \citet{nguyen2018ctdne} \\
\textbf{KDF (integrated)} & \(\exp(-\lambda(C) \cdot dt)\)-form exponential survival decay & \(\deg \le 1\) absolute threshold + \(\theta_U\) sandwich & Laplacian fingerprint analogy & --- \\
\end{longtable}
}

\begin{center}\rule{0.5\linewidth}{0.5pt}\end{center}

\subsection{3.1 Tartaglia et al.~(PNAS 2025) as an Independently Parallel Discovery}\label{tartaglia-et-al.-pnas-2025-as-an-independently-parallel-discovery}

Particularly noteworthy is Tartaglia et al., ``Two-factor synaptic consolidation reconciles robust memory with pruning and homeostatic scaling'' \citep{tartaglia2025twofactor}. \textbf{Its publication overlaps in time with the KDF patent filing (2026-02), and the two works do not cite each other --- exhibiting independent convergence.}

Their claim---``synapse as the product of two factors + replay + homeostatic scaling + Hebbian plasticity → prunes connections while preserving weak memories''---corresponds to our three pillars almost one-to-one (metabolism = homeostatic scaling; rarity protection = two-factor / weak-memory preservation; recombination = Hebbian replay). The fact that two independent teams converged on similar attractor points is \textbf{an observation consistent with the universality hypothesis} (not a proof of the hypothesis).

\begin{center}\rule{0.5\linewidth}{0.5pt}\end{center}

\section{4. Three Structural Similarities of KDF}\label{three-structural-similarities-of-kdf}

\subsection{\texorpdfstring{4.1 \(\delta k^4\) Meta-control ↔ Ginzburg-Landau Quartic Term}{4.1 \textbackslash delta k\^{}4 Meta-control ↔ Ginzburg-Landau Quartic Term}}\label{delta-k4-meta-control-ginzburg-landau-quartic-term}

The Ginzburg-Landau free energy in condensed-matter physics, for an order parameter \(\psi\) near a phase transition, is expanded as

\[
F(\psi) = \alpha_2 \psi^2 + \alpha_4 \psi^4 + \cdots .
\]

\textbf{The quartic term \(\alpha_4 \psi^4\) provides the indispensable restoring force that stabilizes the order parameter at the critical point} (with only the quadratic term, the potential is unbounded below as a square).

KDF's meta-control law \(\Delta \alpha \propto \delta k^4\) matches the Ginzburg-Landau-type potential in \textbf{functional form} on two points: (a) identifying the deviation \(\delta k\) with the order parameter, and (b) generating a quartic restoring force. (This is a functional-form-level correspondence, not a claim of mathematical isomorphism.) Implications:

\begin{itemize}
\tightlist
\item
  KDF's meta-control is implicitly \textbf{steering the network toward self-organization at a critical point}.
\item
  The power-law avalanche distribution predicted by SOC (Self-Organized Criticality) theory may also be expected in KDF (not yet verified; future work).
\item
  Answer to ``why the fourth power'': it is the lowest-order non-trivial restoring force near a phase transition (a physical necessity).
\end{itemize}

\begin{center}\rule{0.5\linewidth}{0.5pt}\end{center}

\subsection{\texorpdfstring{4.2 Sandwich Upper Threshold \(\theta_U\) ↔ Hopfield Spurious-Attractor Rejection}{4.2 Sandwich Upper Threshold \textbackslash theta\_U ↔ Hopfield Spurious-Attractor Rejection}}\label{sandwich-upper-threshold-theta_u-hopfield-spurious-attractor-rejection}

\textbf{Classical problem in Hopfield associative memory.} Storing memory patterns \(\xi^{(1)}, \ldots, \xi^{(P)}\) causes linear combinations of them to emerge as \textbf{spurious attractors}---fixed points that are not themselves learned patterns but behave as false memories. Amit, Gutfreund, and Sompolinsky \citep{amit1987replica}, using the replica analysis of spin-glass models, showed that when the critical capacity \(\alpha_c \equiv P/N \approx 0.138\) is exceeded, spurious attractors proliferate and associative recall collapses. Related phenomena---where repeatedly presented or biased-sampled patterns form broad basins of attraction and obscure other memories, discussed in the context of biased/correlated-pattern capacity---are also closely tied to the spurious problem.

\textbf{No general method for spurious detection has been established in over 40 years.} Modern Hopfield networks \citep{ramsauer2021hopfield} exponentially improved capacity via a softmax mechanism, but a widely accepted general principle for spurious suppression has not yet been presented (this paper does \emph{not} claim that KDF has solved this; see the conjecture at the end of §4.2).

\textbf{KDF's upper threshold \(\theta_U\) formalizes the following heuristic.} When the integrity score \(S\) is ``too perfect'' (for example, \(S > 0.80\)), the candidate is most likely one of:

\begin{enumerate}
\def\labelenumi{\arabic{enumi}.}
\tightlist
\item
  A self-loop or trivial duplicate (rediscovery of an existing neighbor),
\item
  An overfitting / memorization artifact, or
\item
  \textbf{A spurious attractor} (a false fixed point).
\end{enumerate}

By restricting the acceptance region to \([\theta_L, \theta_U]\), these cases are statistically rejected.

\textbf{Mathematical hypothesis.} Integrity scores of random pairs follow a distribution \(p(S)\); true analogies concentrate near the middle of \([\theta_L, \theta_U]\), while spurious attractors concentrate in the right tail \(S > \theta_U\). Upper-bound rejection corresponds to a tail-cut statistical decision boundary.

\textbf{Phase V3 empirical verification (2026-04-18).} On a 100-neuron Hopfield network (Hebbian learning, \(N = 100\), \(P \in \{5, 10, 14, 18, 22\}\)), we implemented and measured a \(\theta\)-filter that rejects the post-recall state as spurious when it has cosine similarity \(\ge \theta\) to multiple learned patterns:

{\def\LTcaptype{none} % do not increment counter
\begin{longtable}[]{@{}
  >{\raggedright\arraybackslash}p{(\linewidth - 6\tabcolsep) * \real{0.2143}}
  >{\centering\arraybackslash}p{(\linewidth - 6\tabcolsep) * \real{0.2143}}
  >{\raggedleft\arraybackslash}p{(\linewidth - 6\tabcolsep) * \real{0.2857}}
  >{\raggedleft\arraybackslash}p{(\linewidth - 6\tabcolsep) * \real{0.2857}}@{}}
\toprule\noalign{}
\begin{minipage}[b]{\linewidth}\raggedright
Detection threshold \(\theta\)
\end{minipage} & \begin{minipage}[b]{\linewidth}\centering
Load factor \(P\)
\end{minipage} & \begin{minipage}[b]{\linewidth}\raggedleft
Spurious-rejection rate
\end{minipage} & \begin{minipage}[b]{\linewidth}\raggedleft
effective\_recall improvement
\end{minipage} \\
\midrule\noalign{}
\endhead
\bottomrule\noalign{}
\endlastfoot
0.80 (KDF canonical) & 18 & \textbf{0\%} & 0 \\
0.70 & 18 & 0\% & 0 \\
0.55 & 18 & 0\% & 0 \\
0.40 & 18 & \textbf{24\%} & 0.49 → 0.65 (+32\%) \\
0.40 & 22 & \textbf{40\%} & 0.34 → 0.56 (+65\%) \\
\end{longtable}
}

\textbf{Findings:}

\begin{itemize}
\tightlist
\item
  {[}True{]} \textbf{The mechanism (multi-pattern similarity rejection) is supported}: an appropriately chosen \(\theta\) can reject \(24\%\) of Hopfield mixture spurious states at load factor \(P = 18\), and up to \(40\%\) at \(P = 22\) (further above the critical capacity \(\alpha_c \approx 0.138\), where spurious attractors proliferate).
\item
  {[}False{]} \textbf{The simple claim of transplanting KDF's canonical \(\theta_U = 0.80\) directly to Hopfield was experimentally refuted}: cosine similarities between a Hopfield mixture state and its constituent patterns each stay around \(\sim 0.4\), undetectable at a \(0.80\) threshold.
\end{itemize}

\textbf{Refinement of the paper's original claim.}

\begin{itemize}
\tightlist
\item
  The original conjecture---``KDF's canonical value \(\theta_U = 0.80\) is the principled-level answer to the Hopfield spurious problem''---\textbf{is refuted}.
\item
  The revised conjecture---``the \textbf{upper-threshold mechanism} is also effective for suppressing Hopfield spurious attractors, but the \textbf{specific value must be tuned domain-specifically}''---is \textbf{partially supported}.
\end{itemize}

Implementation: \href{../../demos/D7_github_issue/src/bin/phase_v3_hopfield_theta_u.rs}{\texttt{demos/D7\_github\_issue/src/bin/phase\_v3\_hopfield\_theta\_u.rs}}. Full record: \href{../VERIFIED_FINDINGS.md}{VERIFIED\_FINDINGS.md F-041}.

Integration verification with modern Hopfield networks \citep{ramsauer2021hopfield} is future work (a minimal experiment for methodology calibration is completed in this phase).

\begin{center}\rule{0.5\linewidth}{0.5pt}\end{center}

\subsubsection{4.2.1 Phase X Step 2 Additional Verification --- Completing the Four-Benchmark Refutation (F-070, 2026-04-19)}\label{phase-x-step-2-additional-verification-completing-the-four-benchmark-refutation-f-070-2026-04-19}

Following the partial refutation of F-041, we systematically re-examined the practical applicability of canonical \((\theta_L, \theta_U) = (0.70, 0.80)\) across three additional benchmarks. Together, they establish four-benchmark cross-cutting evidence.

\textbf{Part A: Sandwich sensitivity analysis on synthetic + F-068 replicated pairs.} We captured the raw scores (with the permissive threshold 0.0) of \texttt{AnalogyDiscoveryEngine::find\_analogy} on 38 pairs in total --- 22 positive (3 Gentner classics including solar system ↔ atom, 4 git ↔ paper cross-domain correspondences, 15 synthetic isomorphic pairs) and 16 negative (1 hand-curated non-isomorphic control, 15 synthetic non-isomorphic pairs) --- and post-hoc applied five sandwich conditions \(\{(0.70, 0.75), (0.70, 0.80), (0.70, 0.90), (0.70, 0.95), (0.70, 1.00)\}\). Results:

{\def\LTcaptype{none} % do not increment counter
\begin{longtable}[]{@{}
  >{\raggedright\arraybackslash}p{(\linewidth - 10\tabcolsep) * \real{0.1304}}
  >{\raggedleft\arraybackslash}p{(\linewidth - 10\tabcolsep) * \real{0.1739}}
  >{\raggedleft\arraybackslash}p{(\linewidth - 10\tabcolsep) * \real{0.1739}}
  >{\raggedleft\arraybackslash}p{(\linewidth - 10\tabcolsep) * \real{0.1739}}
  >{\raggedleft\arraybackslash}p{(\linewidth - 10\tabcolsep) * \real{0.1739}}
  >{\raggedleft\arraybackslash}p{(\linewidth - 10\tabcolsep) * \real{0.1739}}@{}}
\toprule\noalign{}
\begin{minipage}[b]{\linewidth}\raggedright
\((\theta_L, \theta_U)\)
\end{minipage} & \begin{minipage}[b]{\linewidth}\raggedleft
TP
\end{minipage} & \begin{minipage}[b]{\linewidth}\raggedleft
FN
\end{minipage} & \begin{minipage}[b]{\linewidth}\raggedleft
TN
\end{minipage} & \begin{minipage}[b]{\linewidth}\raggedleft
FP
\end{minipage} & \begin{minipage}[b]{\linewidth}\raggedleft
F1
\end{minipage} \\
\midrule\noalign{}
\endhead
\bottomrule\noalign{}
\endlastfoot
\((0.70, 0.75)\) & 0 & 22 & 16 & 0 & 0.000 \\
\textbf{\((0.70, 0.80)\) canonical} & \textbf{0} & \textbf{22} & \textbf{16} & \textbf{0} & \textbf{0.000} \\
\((0.70, 0.90)\) & 0 & 22 & 16 & 0 & 0.000 \\
\((0.70, 0.95)\) & 0 & 22 & 16 & 0 & 0.000 \\
\((0.70, 1.00)\) & 22 & 0 & 16 & 0 & \textbf{1.000} \\
\end{longtable}
}

\textbf{Finding.} Positive-analogy scores saturate uniformly at \(\ge 0.99\) (fingerprint identity for graph isomorphism), while negatives distribute around \(\sim 0.57\)--\(0.60\). Because canonical \(\theta_U = 0.80\) rejects every positive, F1 is 0.000; conversely, \(\theta_U = 1.00\) (no effective upper bound) yields complete separation via \(\theta_L = 0.70\) alone and reaches F1 = 1.000. \textbf{No empirical scores fall within the ``middle band'' that 0.80 would admit.}

\textbf{Part B: KdfProcessorRev12 review loop on LoCoMo temporal 30 Q.} Fixing \((t_{\text{wait1}}, t_{\text{wait2}}) = (30, 30)\), we ran the review cycle under three conditions \(\theta_U \in \{0.80, 0.90, 1.00\}\) up to 65 cycles. The spoke\_up / demote distribution across the total 1,140 RARE nodes (of which 8 belong to answer turns):

{\def\LTcaptype{none} % do not increment counter
\begin{longtable}[]{@{}
  >{\raggedright\arraybackslash}p{(\linewidth - 14\tabcolsep) * \real{0.0968}}
  >{\raggedleft\arraybackslash}p{(\linewidth - 14\tabcolsep) * \real{0.1290}}
  >{\raggedleft\arraybackslash}p{(\linewidth - 14\tabcolsep) * \real{0.1290}}
  >{\raggedleft\arraybackslash}p{(\linewidth - 14\tabcolsep) * \real{0.1290}}
  >{\raggedleft\arraybackslash}p{(\linewidth - 14\tabcolsep) * \real{0.1290}}
  >{\raggedleft\arraybackslash}p{(\linewidth - 14\tabcolsep) * \real{0.1290}}
  >{\raggedleft\arraybackslash}p{(\linewidth - 14\tabcolsep) * \real{0.1290}}
  >{\raggedleft\arraybackslash}p{(\linewidth - 14\tabcolsep) * \real{0.1290}}@{}}
\toprule\noalign{}
\begin{minipage}[b]{\linewidth}\raggedright
\(\theta_U\)
\end{minipage} & \begin{minipage}[b]{\linewidth}\raggedleft
total RARE
\end{minipage} & \begin{minipage}[b]{\linewidth}\raggedleft
answer-RARE
\end{minipage} & \begin{minipage}[b]{\linewidth}\raggedleft
spoke\_up (ans)
\end{minipage} & \begin{minipage}[b]{\linewidth}\raggedleft
demoted (ans)
\end{minipage} & \begin{minipage}[b]{\linewidth}\raggedleft
spoke\_up (non-ans)
\end{minipage} & \begin{minipage}[b]{\linewidth}\raggedleft
demoted (non-ans)
\end{minipage} & \begin{minipage}[b]{\linewidth}\raggedleft
avg cycles
\end{minipage} \\
\midrule\noalign{}
\endhead
\bottomrule\noalign{}
\endlastfoot
\textbf{0.80 canonical} & 1140 & 8 & \textbf{0} & \textbf{8 (100\%)} & 0 & \textbf{1132 (100\%)} & 60.0 \\
0.90 & 1140 & 8 & 8 (100\%) & 0 & 1132 (100\%) & 0 & 1.0 \\
1.00 & 1140 & 8 & 8 (100\%) & 0 & 1132 (100\%) & 0 & 1.0 \\
\end{longtable}
}

\textbf{Finding.} Under canonical \(\theta_U = 0.80\), every RARE node in the LoCoMo chain graph is demoted to Garbage after the 60-cycle timeout (complete information loss). Under \(\theta_U \ge 0.90\), all RARE nodes spoke up within one cycle, but there is no discrimination between answer and non-answer turns, so the filter does not function. \textbf{The sparse chain structure of LoCoMo cannot admit a discriminative middle band under the sandwich.}

\textbf{Cross-benchmark verdict:}

{\def\LTcaptype{none} % do not increment counter
\begin{longtable}[]{@{}
  >{\raggedright\arraybackslash}p{(\linewidth - 4\tabcolsep) * \real{0.3333}}
  >{\raggedright\arraybackslash}p{(\linewidth - 4\tabcolsep) * \real{0.3333}}
  >{\raggedright\arraybackslash}p{(\linewidth - 4\tabcolsep) * \real{0.3333}}@{}}
\toprule\noalign{}
\begin{minipage}[b]{\linewidth}\raggedright
Benchmark
\end{minipage} & \begin{minipage}[b]{\linewidth}\raggedright
Domain
\end{minipage} & \begin{minipage}[b]{\linewidth}\raggedright
Verdict on canonical \((0.70, 0.80)\)
\end{minipage} \\
\midrule\noalign{}
\endhead
\bottomrule\noalign{}
\endlastfoot
F-041 Hopfield mixture & Associative memory & 0\% detection at \(\theta_U = 0.80\); \(24\%\) at \(P = 18\) and \(40\%\) at \(P = 22\) only after lowering to \(\theta = 0.40\) \\
F-068 analogy-engine direct & Graph isomorphism & Scores at \(\ge 0.99\); canonical rejects every true positive \\
F-070 Part A (38 pairs) & Synthetic + F-068 replication & F1(canonical) \(= 0.000\); F1\(((0.70, 1.00)) = 1.000\) \\
F-070 Part B (LoCoMo streaming) & 30 Q × 1,140 RARE & 100\% RARE demoted under canonical → complete information loss \\
\end{longtable}
}

→ \textbf{Across four benchmarks, canonical \((\theta_L, \theta_U) = (0.70, 0.80)\) loses practical value.} True positive-analogy scores concentrate at \(\ge 0.95\) and negatives at \(\le 0.65\), so the empirically correct sandwich is \((0.70, 1.00)\) (effectively no upper bound; only \(\theta_L\) active) or at the \((0.90, 1.00)\) level.

\textbf{Final narrowing of the claim.}

\begin{itemize}
\tightlist
\item
  {[}False{]} Original conjecture --- ``canonical \((0.70, 0.80)\) is the principled-level answer to Hopfield spurious / over-similar candidates'' --- \textbf{is refuted across four benchmarks}.
\item
  {[}True{]} Novelty of the 2-threshold \emph{mechanism} (the structure that admits only the middle band by simultaneously imposing lower and upper bounds) has no counterpart among the ten related disciplines we surveyed; the novelty claim is \textbf{maintained at the mechanism level}.
\item
  {[}Calibration Required{]} Specific canonical values require domain-specific empirical calibration. The score distributions from F-068 / F-070 suggest \(\theta_U \ge 0.95\) for graph-isomorphism regimes and \(\theta \le 0.5\) for associative-memory regimes.
\end{itemize}

Implementation: \href{../../demos/D8_llm_memory/src/bin/phase_x2_sandwich_twait_locomo.rs}{\texttt{demos/D8\_llm\_memory/src/bin/phase\_x2\_sandwich\_twait\_locomo.rs}}. Full record: \href{../VERIFIED_FINDINGS.md}{VERIFIED\_FINDINGS.md F-070}.

\begin{center}\rule{0.5\linewidth}{0.5pt}\end{center}

\subsection{\texorpdfstring{4.3 \(\exp(-\lambda(C) \cdot dt)\) ↔ Locally-Varying Continuous-Time Markov Chain}{4.3 \textbackslash exp(-\textbackslash lambda(C) \textbackslash cdot dt) ↔ Locally-Varying Continuous-Time Markov Chain}}\label{exp-lambdac-cdot-dt-locally-varying-continuous-time-markov-chain}

The survival probability of a continuous-time Markov process \(X_t\) (with exponentially distributed residence times) is generally written as

\[
\Pr[\text{no jump in } [t, t+dt]] = \exp(-\Lambda(x) \cdot dt) .
\]

KDF's weight update \(w \leftarrow w \cdot \exp(-\lambda(C) \cdot dt)\) has \textbf{the same functional form}. If we interpret \(\lambda(C)\) as a local ``jump rate'', KDF's edge decay corresponds to the \textbf{motif of a survival process}---``the probability that an edge is lost in an infinitesimal interval follows \(1 - \exp(-\lambda(C) \cdot dt)\)''.

\textbf{However, strict CTMC equivalence does not hold.} In KDF, (a) \(\lambda(C)\) changes as adjacent edges are pruned, so the generator is non-stationary; (b) the weight \(w\) is a continuous quantity, not a probability mass; (c) edge protection (Rare) depends on global conditions; and so on --- the standard CTMC assumptions are not satisfied. This section's claim therefore remains at the level of \textbf{motivation for a theoretical connection (a motivating analogy)}; ergodic results such as stationary distributions or spectral gaps cannot be imported directly.

What this structural correspondence suggests is the following --- \textbf{directions for future theoretical development}, not theorems:

\begin{itemize}
\tightlist
\item
  Claim 46's Laplacian eigenvalue fingerprint is \textbf{conceptually consistent} with the attempt to project ego-graph spectral-gap information into a 32-dimensional space (an exact information-preservation proof has not been carried out).
\item
  Theoretical connections to diffusion maps, PageRank, and spectral clustering are future work.
\end{itemize}

\begin{center}\rule{0.5\linewidth}{0.5pt}\end{center}

\section{5. Empirical Evaluation}\label{empirical-evaluation}

\subsection{5.1 Positive Results (Six Cases)}\label{positive-results-six-cases}

{\def\LTcaptype{none} % do not increment counter
\begin{longtable}[]{@{}
  >{\raggedright\arraybackslash}p{(\linewidth - 6\tabcolsep) * \real{0.2500}}
  >{\raggedright\arraybackslash}p{(\linewidth - 6\tabcolsep) * \real{0.2500}}
  >{\raggedright\arraybackslash}p{(\linewidth - 6\tabcolsep) * \real{0.2500}}
  >{\raggedright\arraybackslash}p{(\linewidth - 6\tabcolsep) * \real{0.2500}}@{}}
\toprule\noalign{}
\begin{minipage}[b]{\linewidth}\raggedright
Problem
\end{minipage} & \begin{minipage}[b]{\linewidth}\raggedright
Data
\end{minipage} & \begin{minipage}[b]{\linewidth}\raggedright
KDF result
\end{minipage} & \begin{minipage}[b]{\linewidth}\raggedright
Comparison
\end{minipage} \\
\midrule\noalign{}
\endhead
\bottomrule\noalign{}
\endlastfoot
P1 LLM-agent memory & LongMemEval 500 Q (ICLR 2025) & Recall \(= 0.821\) & Industry-default TTL \(= 0.107\) (\textbf{\(\times 7.7\)}); Random \(= 0.294\) (\(\times 2.8\)) \\
P2 Personal knowledge base & Obsidian vault of 2,182 notes (PII-masked) & \(F_1 = 0.747\) & Wilcoxon \(p = 0.006\) vs.~Random / OrphanOnly / TextSim \\
P3 Large-scale log observability (static baseline) & NASA HTTP log, 50k real records & Recall \(= 0.237\) (keep \(10\%\), \(4xx/5xx\) retention) & Random \(= 0.102\) (\textbf{\(\times 2.3\)}), without labels \\
P7 ML reproducibility meta (by-product) & 5 benchmarks × 5 scenarios & 4/5 exact predictions + 1 via alternate route & Released independently as the \texttt{bias-detector} crate \\
\textbf{P8 Distributed-execution bit-exactness (Claim 17)} & \textbf{LoCoMo 10 real graph (600+ nodes, 400+ edges)} & \textbf{max edge-weight diff \(= 0.000 \text{e}0\)} & \texttt{apply\_edge\_decay} (global) and \texttt{apply\_edge\_decay\_local} (distributed) are completely bit-exact, F-069 \\
\textbf{P11 NASA streaming: Claim 14 decay benefit} & \textbf{NASA HTTP log 50k real records, replayed in time order (500 rec / 100 window)} & \textbf{C1 decay rare\_recall \(= 0.4898\) at keep \(30\%\) (C0 static \(0.4592\) → \(+3.06\) pt)} & \textbf{First empirical anchor for the streaming-benefit thesis {[}v2: scoped to \emph{temporally recurring rare} per F-087; see §5.x Phase 2/2.5 subsection{]} {[}v2.1: F-093 で本 dataset の rare = 4xx/5xx は実質 404-pattern driven、5xx response は本 dataset に不在のため generalization 未測定{]} {[}v2.2: F-094 で Apache 別 dataset を recurring rare 定義 (freq \(\ge 5\)) で再 test、\(+3.67\) pt 再現、cross-domain N=2 で recurring-rare 軸が durable{]} {[}v2.5: F-097 で BGL HW kernel log を recurring rare で再 test、\(+33.33\) pt 再現、cross-domain N=3 (web + HW family) で軸 durable、ただし small-alphabet domain (8 templates) では sanity V\_one-shot inconsistent{]} {[}v2.7: F-100 で HDFS distributed file system log を H\_R+\_literal で test、\(-23.08\) pt FAIL、H\_anomaly も \(-4.30\) pt FAIL、N=4 cross-family 不能、durability arc を web + HW kernel log family specific に narrow{]}, F-072} \\
\end{longtable}
}

As additional verification in Phase X Step 1, all three mechanisms of Claim 1---metabolic control, rarity protection, and integrity discovery---are now empirically backed on realistic benchmarks (integrity discovery via F-068: \(100\%\) on Gentner classics, \(100\%\) on cross-domain git↔paper, \(0\%\) false positives on the negative control). In Phase X Step 5 (F-072), Claim 14 exponential decay is confirmed to yield a \(+3.06\)-point benefit in a realistic streaming scenario, giving \textbf{empirical anchor to the narrowed thesis that ``streaming is the true use case''}. \textbf{{[}v2{]}} F-087 Apache error log replication (\(-13.04\) pt vs.~C0 static, sign reversal) further restricts this anchor to \textbf{temporally recurring rare events} (e.g., the persistent failure pattern of NASA HTTP 4xx/5xx). For one-shot rare events (e.g., reconnaissance probes), streaming with decay is actively harmful --- see §5 Phase 2/2.5 subsection.

\begin{center}\rule{0.5\linewidth}{0.5pt}\end{center}

\subsection{5.2 Negative Results (Four Cases) --- An Honest Record}\label{negative-results-four-cases-an-honest-record}

For the integrity of this work, we report \textbf{both generalization failures found in testing} and \textbf{refutations of canonical values that we ourselves proposed}:

{\def\LTcaptype{none} % do not increment counter
\begin{longtable}[]{@{}
  >{\raggedright\arraybackslash}p{(\linewidth - 4\tabcolsep) * \real{0.3333}}
  >{\raggedright\arraybackslash}p{(\linewidth - 4\tabcolsep) * \real{0.3333}}
  >{\raggedright\arraybackslash}p{(\linewidth - 4\tabcolsep) * \real{0.3333}}@{}}
\toprule\noalign{}
\begin{minipage}[b]{\linewidth}\raggedright
Problem
\end{minipage} & \begin{minipage}[b]{\linewidth}\raggedright
Result
\end{minipage} & \begin{minipage}[b]{\linewidth}\raggedright
Interpretation
\end{minipage} \\
\midrule\noalign{}
\endhead
\bottomrule\noalign{}
\endlastfoot
P6 OSS maintenance & KDF/Random ratios over three repos (\texttt{rust-lang/rust}, \texttt{tokio-rs/tokio}, \texttt{golang/go}) are \(\times 1.13, \times 1.03, \times 0.85\); average \textbf{\(\times 1.00\)} & The \(+15\%\) on \texttt{rust-lang} alone is a repository-specific local signal. The general-OSS applicability claim is \textbf{withdrawn} (F-038) \\
P5 Paper rediscovery & OpenAlex 200 papers × concept-sharing graph; KDF/Random \(= \mathbf{\times 0.83}\) & Late-bloomer detection is a D5-type task (independent of structure); KDF is ill-suited to concept-graphs (F-039) \\
\textbf{P9 Redundancy of Claim 5/14 time signals on static tasks (validated on streaming)} & \textbf{LoCoMo 321 Q at keep \(30\%\): KDF\_static \(= 0.5286\); adding time signals yields \(0.43\)--\(0.53\) (static: all worse or tied). NASA streaming 50k records at keep \(30\%\): C0 static \(0.4592\); C1 Claim 14 decay \(0.4898\) (}\(+3.06\) pt benefit on streaming\textbf{)} & On static query tasks, structural rarity already subsumes temporal rarity, so the time signals are redundant. On streaming, decay discards stale normal traffic and relatively elevates rare resources → \textbf{task-structure-dependent conditional value} (F-069 static redundancy + F-072 streaming \(+3.06\)-pt validation) \\
\textbf{P10 Four-benchmark refutation of canonical \(\theta_U = 0.80\) (Claims 47--48)} & \textbf{F-041 (Hopfield: \(0\%\) detection) + F-068 (analogy scores \(0.99+\) reject every positive) + F-070 Part A (F1 \(0.000\) vs.~\(1.000\)) + F-070 Part B (LoCoMo: \(100\%\) RARE demoted)} & The 2-threshold \emph{mechanism} is supported; the canonical specific values \((0.70, 0.80)\) are refuted across four benchmarks. Domain-specific calibration is required (F-041, F-070) \\
\end{longtable}
}

\begin{center}\rule{0.5\linewidth}{0.5pt}\end{center}

\subsection{\texorpdfstring{5.3 A Priori Applicability Metric: \texttt{bias-detector}}{5.3 A Priori Applicability Metric: bias-detector}}\label{a-priori-applicability-metric-bias-detector}

We released a zero-dependency Rust crate for determining KDF's applicability in advance (\href{../../crates/bias-detector/}{\texttt{crates/bias-detector/}}). The metric \(\text{bias\_score} = 0.3 \cdot I_1 + 0.7 \cdot I_4\) correctly predicted applicability on 4 of 5 benchmarks exactly, and on the 5th via an alternate route.

\begin{center}\rule{0.5\linewidth}{0.5pt}\end{center}

\section{6. Discussion / Implications}\label{discussion-implications}

\subsection{6.1 The Possibility of a Domain-Invariant Architecture}\label{the-possibility-of-a-domain-invariant-architecture}

From the qualitative observation that ten independent disciplines adopt the same three-pillar structure (or a 2--3-pillar subset), we advance the following \textbf{hypothesis} (not a proof):

\begin{quote}
\textbf{Hypothesis (untested)}: ``Systems that preserve information long-term under finite resources repeatedly exhibit designs belonging to the architectural family of three mechanisms---metabolism, rarity protection, and recombination.''
\end{quote}

The hypothesis \textbf{can be supported} (a necessary but not sufficient set of grounds) from the following directions:

\begin{enumerate}
\def\labelenumi{(\alph{enumi})}
\tightlist
\item
  independent evolutionary convergence in biological systems (brain, immune);
\item
  independent rediscovery in ML engineering (EWC, ECS);
\item
  functional-form match with critical-phenomena physics.
\end{enumerate}

However, all three are observations of \emph{compatibility}, not proofs of \emph{necessity}. Demonstrating genuinely universal optimality would require the following---none carried out in this paper:

\begin{itemize}
\tightlist
\item
  Mathematical formalization of the finite-resource information-preservation problem and characterization of its optimal-solution class.
\item
  Proof of a performance lower bound on this problem for systems that \textbf{lack} the three pillars.
\item
  Statistical tests quantifying the ten-domain correspondence.
\end{itemize}

KDF in this work \textbf{provides the first comprehensive numerical implementation} in the direction of this hypothesis; verifying the hypothesis itself is a task for future work.

\begin{center}\rule{0.5\linewidth}{0.5pt}\end{center}

\subsection{6.2 The Possibility of Domain-Specialized Implementations}\label{the-possibility-of-domain-specialized-implementations}

By wrapping the same core engine with domain-specific interfaces, we can extend deployment across multiple application markets:

\begin{itemize}
\tightlist
\item
  \texttt{kdf-associative-memory} --- a wrapper for Hopfield spurious-attractor suppression
\item
  \texttt{kdf-coreset} --- unsupervised, label-free Equitable Coreset Selection
\item
  \texttt{kdf-temporal-graph} --- temporal-graph embedding, directly compared with \citet{nguyen2018ctdne}
\item
  \texttt{kdf-portfolio} --- information tail-risk management (insurance / finance applications)
\item
  \texttt{kdf-llm-memory} --- LLM-agent long-term memory (leveraging the LongMemEval track record directly)
\end{itemize}

\begin{center}\rule{0.5\linewidth}{0.5pt}\end{center}

\subsection{6.3 Positioning of Patents and Licensing}\label{positioning-of-patents-and-licensing}

This paper's patent claims secure two strategic contributions:

\begin{enumerate}
\def\labelenumi{\arabic{enumi}.}
\tightlist
\item
  \textbf{Claim 1 (independent)}: integration of the three mechanisms. Even if the individual elements pre-exist, the integration itself has no prior art. F-068 completed the realistic benchmark for the analogy mechanism, so all three mechanisms are now empirically backed.
\item
  \textbf{Claims 47--48: the sandwich 2-threshold \emph{mechanism}} (lower bound \(\theta_L\) + upper bound \(\theta_U\)) has no counterpart in the ten related disciplines surveyed and is a distinctive element. However, the patent-designated canonical values \((\theta_L, \theta_U) = (0.70, 0.80)\) themselves were subject to the four-benchmark refutation in §4.2 / §5.2 P10, narrowing the claim to \emph{``the mechanism is novel; specific values require domain-calibration.''}
\end{enumerate}

The implementation code is available under PolyForm Noncommercial 1.0.0 (research, education, and personal use freely permitted). Commercial licenses --- both for the source code and for practicing the patent --- are managed separately; inquiries are welcome via the repository's COMMERCIAL.md.

\begin{center}\rule{0.5\linewidth}{0.5pt}\end{center}

\subsection{6.4 Limitations}\label{limitations}

\begin{itemize}
\tightlist
\item
  \textbf{The ``domain-invariant architecture'' hypothesis itself is untested.} The ten-domain correspondence is merely a qualitative observation; a formal theorem of ``necessary convergence'' has not been proved. §6.1's hypothesis should be read \textbf{as a hypothesis, not a conclusion}.
\item
  \textbf{Canonical values \((\theta_L, \theta_U) = (0.70, 0.80)\) are refuted across four benchmarks} (F-041 Hopfield / F-068 analogy / F-070 Part A synthetic / F-070 Part B LoCoMo streaming). The sandwich mechanism itself is supported, but universality of specific values is not claimed (§4.2 / §5.2 P10).
\item
  \textbf{Claim 5 / 14 time-evaluation components and exponential decay are task-structure-dependent}: redundant on static query tasks (F-069 LoCoMo); \(+3.06\)-pt benefit on streaming scenarios (F-072 NASA 50 k records, time-ordered replay). When structural rarity already subsumes temporal rarity, the time signals are redundant; when stale traffic must be discarded under continuous operation, they are effective. \textbf{{[}v2{]}} F-087 Apache error log replication (\(-13.04\) pt) further narrows the streaming benefit to \emph{temporally recurring rare events} (persistent failure modes); for \emph{one-shot rare events} (e.g., reconnaissance probes), decay actively erodes recall and ActivationScore amplifies common-path dominance. \textbf{{[}v2.5{]}} F-097 BGL HW kernel log replication (\(+33.33\) pt on recurring-rare) extends the durable axis to N=3 cross-family (NASA HTTP + Apache error + BGL HW kernel), but the sanity inconsistency (\(+0\) pt on V\_one-shot vs.~F-087 / F-094 negative direction) narrows the \emph{one-shot disposal mechanism} to web log family with rich resource alphabet --- small-alphabet domains lack streaming discrimination room. \textbf{{[}v2.7{]}} F-100 HDFS distributed file system log replication FAILED both hypotheses (H\_R+\_literal \(\Delta = -23.08\) pt FAIL, H\_anomaly \(\Delta = -4.30\) pt FAIL): cross-domain durability of the recurring-rare benefit is narrowed to \textbf{web access log family + HW kernel log family} specific; \textbf{distributed file system log domain is out of scope} of current empirical claims, and N=4 cross-family expansion is not achieved.
\item
  \textbf{Claim 25 ActivationScore depends on the temporal distribution of rare events}: \(100\%\) rescue for temporally clustered rare events (F-027 Mode E drift scenario); hurts or neutralizes on evenly distributed rare events (F-072 NASA; F-069 LoCoMo) due to recency bias. At application time, discrimination of the rare-event temporal pattern is required. \textbf{{[}v2{]}} F-087 Apache error log replication quantifies this further: with one-shot rare items (freq \(\le 10\) resource paths), full streaming (decay + activation + meta α) reaches recall \(0.000\) vs.~C0 static \(0.4348\) (\(-43.48\) pt). ActivationScore \emph{inverts} the desired ordering when rare items lack temporal recurrence.
\item
  \textbf{The ten-domain correspondences are at the level of structural similarity.} Proof as mathematical isomorphism (structure-preserving bijection) has not been carried out for any of the disciplines.
\item
  \textbf{The §4.3 Markov correspondence is a motivating analogy.} Strict CTMC equivalence does not hold (non-stationary generator, weights are not probability mass, etc.).
\item
  \textbf{Empirical evidence covers only five domains with six cases} (P1 / P2 / P3 / P7 / P8 = Claim 17 bit-exact; P11 = NASA streaming; P3 and P11 are the static and streaming scenarios of the same NASA HTTP-log domain). Generalization claims beyond these five domains are \textbf{beyond the scope of this paper}.
\item
  \textbf{Negative results P5 / P6} establish that KDF is ineffective for D5-type tasks (rarity independent of structure). \textbf{P9 / P10} establish that our own canonical parameters are not optimal in realistic scenarios. A priori screening via \texttt{bias-detector} (F-030 / F-036) is recommended.
\item
  \textbf{Optimality of specific parameter values} (\(\beta = 0.01\), \(7:2:1\), \(5:3:1\), \([0.70, 0.80]\), \(\delta k^4\), \(\lambda\), \(\tau_{\text{ref}}\)) is empirically chosen; domain-specific tuning is required. P10 quantifies one such case, where the sandwich canonical fails to match the true score distribution.
\item
  \textbf{Universality \(\neq\) optimality \(\neq\) novelty.} The ten-discipline parallel discovery suggests universality, but whether the common structure is \textbf{genuinely optimal} is a separate question; and novelty is limited to the narrow three-item claim in §1.4 (novelty of integration + novelty of mechanism + novelty of open implementation).
\end{itemize}

\begin{center}\rule{0.5\linewidth}{0.5pt}\end{center}

\section{7. Conclusion}\label{conclusion}

We presented the Knowledge Decay Framework (KDF) as a three-mechanism integrated architecture for operating information networks under finite resources. The reference implementation \texttt{cgb-kdf}, corresponding to claims 1--50 of patent JP 2026-027032, carries direct verification tests \texttt{test\_claimN\_*} for all 50 claims (56 tests total); the workspace excluding \texttt{kdf-python} / \texttt{kdf-wasm} passes 449 tests (measured 2026-04-18; see \href{../patent/COMPLIANCE.md}{\texttt{COMPLIANCE.md}}).

Our central contributions are three:

\begin{enumerate}
\def\labelenumi{\arabic{enumi}.}
\tightlist
\item
  \textbf{A qualitative systematization of the three-pillar structure across ten independent disciplines}, with KDF advanced as \textbf{a candidate implementation} of a domain-invariant architectural family (formal proof of ``necessary convergence'' not carried out). All three mechanisms of Claim 1 (metabolic control / rarity protection / integrity discovery) reach empirically backed status on realistic benchmarks via F-068 / F-069 / F-070.
\item
  \textbf{Three structural similarities} (\(\delta k^4\) ↔ Ginzburg-Landau quartic term {[}functional-form match{]}; sandwich 2-threshold \emph{mechanism} ↔ Hopfield spurious rejection {[}mechanism supported, canonical values refuted in §4.2 on four benchmarks{]}; \(\exp(-\lambda \cdot dt)\) ↔ Markov survival probability {[}motivating analogy{]}) indicate \textbf{directions of theoretical connection} for KDF's engineering heuristics (rigorous theorem formulation is future work).
\item
  \textbf{Honest empirical reporting} across three streams: positive results (P1 / P2 / P3 / P7 / P8 = Claim 17 bit-exact; P11 = NASA streaming \(+3.06\) pt; six total); generalization failures on comparison benchmarks (P5 / P6); and \textbf{refutations of canonical values we proposed ourselves (P9 Claim 5/14 static-task redundancy; P10 Claim 47--48 four-benchmark refutation of canonical values)} (four negatives total). We also provide the \texttt{bias-detector} a priori applicability metric. Self-refuting our own claim across four benchmarks serves the dual purpose of novelty narrowing and credibility strengthening.
\end{enumerate}

KDF is \textbf{not a universal solution}. It cannot directly contribute to physical problems (climate change), economic problems (poverty), or social problems (war, educational inequality). \textbf{Moreover, the BEIR SciFact verification on 2026-04-18 (F-045) empirically confirmed that KDF is not applicable to general semantic retrieval (query--document matching)} (recall@10 \(= 0.000\), below Random).

\textbf{The refined applicability predictor, established with the additional verifications of 2026-04-19 (F-061--F-065), is:}

\begin{quote}
\textbf{``KDF decisively outperforms Random and baseline heuristics only under the condition that structural rareness correlates with task importance.''}
\end{quote}

Task type × KDF applicability matrix:

{\def\LTcaptype{none} % do not increment counter
\begin{longtable}[]{@{}
  >{\raggedright\arraybackslash}p{(\linewidth - 6\tabcolsep) * \real{0.2500}}
  >{\raggedright\arraybackslash}p{(\linewidth - 6\tabcolsep) * \real{0.2500}}
  >{\centering\arraybackslash}p{(\linewidth - 6\tabcolsep) * \real{0.2500}}
  >{\raggedright\arraybackslash}p{(\linewidth - 6\tabcolsep) * \real{0.2500}}@{}}
\toprule\noalign{}
\begin{minipage}[b]{\linewidth}\raggedright
Task
\end{minipage} & \begin{minipage}[b]{\linewidth}\raggedright
structural signal vs.~importance
\end{minipage} & \begin{minipage}[b]{\linewidth}\centering
KDF applicability
\end{minipage} & \begin{minipage}[b]{\linewidth}\raggedright
Evidence
\end{minipage} \\
\midrule\noalign{}
\endhead
\bottomrule\noalign{}
\endlastfoot
Path-based algorithms (APSP) & path-critical = bottleneck = rare & {[}True{]} & F-061 \\
Integration-point preservation (git merges, \textbf{repo merge rate \(< 10\%\)}) & merge = high degree = important & {[}True{]} & F-062, F-065 \\
LLM-memory temporal recall (long-context date/time) & rare date/time literal = structurally rare & {[}True{]} & F-057 / F-058 \\
Orphan-note detection (PKM) & orphan = deg \(0\) = rare & {[}True{]} & F-012 / F-017 \\
\textbf{Merge-heavy repos (merge rate \(> 20\%\))} & \textbf{merge \(\neq\) rare; TopDegree wins} & {[}False{]} & F-065 pytest \\
\textbf{GP / kernel-regression inducing points} & \textbf{density center \(\neq\) rareness} & {[}False{]} & F-063 \\
\textbf{Naive Python call-graph API} & \textbf{API = high in-degree; opposite of KDF Rare} & {[}False{]} & F-064 \\
Scale-free hub centrality & degree ≈ betweenness; TopDegree suffices & {[}False{]} & F-061 BA/WS \\
Metadata-based minority (cultural / semantic) & metadata independent of structure & {[}False{]} & F-047 \\
General semantic retrieval & semantic understanding is essential & {[}False{]} & F-045 \\
\end{longtable}
}

KDF is effective only under the following restricted conditions:

\begin{itemize}
\tightlist
\item
  {[}True{]} the data can be represented as a graph;
\item
  {[}True{]} \textbf{the query is unknown at retention time} (to be searched for later);
\item
  {[}True{]} \textbf{protection of one-off / rarely-mentioned items} is the goal;
\item
  {[}True{]} most of the data is redundant, a minority is unique, and labeling is difficult;
\item
  {[}True{]} \textbf{task importance correlates with ``structural rareness''} (newly refined via F-061--F-065).
\end{itemize}

Under these constraints, KDF is effective for \textbf{memory-curation tasks such as conversational memory and PKM} (P1: LongMemEval TTL \(\times 7.7\); P2: Obsidian F1 \(= 0.747\), \(p = 0.006\); P3: NASA log \(\times 2.3\)).

\textbf{The initial F-044 measurement on 2026-04-18 turned out to be a simulation (the ``KDF'' numbers used a constant 0.821 recall value inherited from F-033, an earlier static LongMemEval measurement, rather than running the real KDF algorithm), as determined by F-052 / F-053; re-running with the real KDF completely reversed the conclusion on LongMemEval and yielded partial wins on LoCoMo:}

\begin{itemize}
\tightlist
\item
  \textbf{LongMemEval 500 Q (F-053)}: Real KDF \(0.434\) vs.~Mem0 \(0.672\) (Mem0 \(+23.8\) pt, \(p < 10^{-16}\)). Mem0 wins significantly in 5 of 6 categories.
\item
  \textbf{LoCoMo 200 Q balanced (F-056)}: Real KDF \(0.535\) vs.~Mem0 \(0.590\) (gap \(-5.5\) pt, \textbf{\(p = 0.24\), statistically tied}).

  \begin{itemize}
  \tightlist
  \item
    \textbf{Temporal category (LoCoMo 50 Q sub)}: KDF \(0.460\) vs.~Mem0 \(0.240\) (\textbf{KDF \(+22\) pt, \(p = 0.035\) win}).
  \item
    Mem0 wins narrative by \(+24\) pt; factual / inferential are n.s.
  \end{itemize}
\item
  \textbf{LoCoMo temporal full 321 Q, gpt-4o-mini (F-057)}: Real KDF \(0.312\) vs.~Mem0 \(0.206\) (\textbf{KDF \(+10.6\) pt, \(p = 0.0014\)}, McNemar \(b = 71\) / \(c = 37\)).
\item
  \textbf{LoCoMo temporal full 321 Q, gpt-4.1-mini (F-058, model-robustness check)}: Real KDF \(0.324\) vs.~Mem0 \(0.090\) (\textbf{KDF \(+23.4\) pt, \(p = 1.6 \times 10^{-14}\)}, McNemar \(b = 89\) / \(c = 14\)).

  \begin{itemize}
  \tightlist
  \item
    \textbf{Updating to the newer-generation model degrades Mem0 while KDF is preserved, widening the gap.}
  \item
    Hypothesis: gpt-4.1-mini's fact extraction performs more aggressive compression and loses even more temporal information.
  \item
    KDF's temporal advantage is \textbf{robust across 2 models × 321 Q, exposing a principled weakness of fact-extraction-based memory systems}.
  \end{itemize}
\item
  \textbf{LongMemEval 500 Q × gpt-4.1-mini (F-059, completing the 2×2 matrix)}: Real KDF \(0.452\) vs.~Mem0 \(0.722\) (\textbf{Mem0 \(+27.0\) pt, \(p = 3.06 \times 10^{-23}\)}, McNemar \(b = 32\) / \(c = 167\)).

  \begin{itemize}
  \tightlist
  \item
    \textbf{The F-053 gap of \(-23.8\) pt widens to \(-27.0\) pt} (Mem0 strengthens on the new model; KDF marginally increases).
  \item
    Per-category: Mem0 \(+61\)-pt landslide on single-session-assistant; temporal-reasoning gap \(-6\) pt, \(p = 0.23\) n.s. (KDF remains competitive on short-dialogue temporal queries).
  \item
    → \textbf{LongMemEval is a benchmark that showcases Mem0's structural strength}; KDF's room to win is limited to n.s. categories (temporal-reasoning, single-session-preference).
  \end{itemize}
\end{itemize}

\textbf{Completed 2 benchmarks × 2 models matrix (F-053 / F-057 / F-058 / F-059)}:

{\def\LTcaptype{none} % do not increment counter
\begin{longtable}[]{@{}
  >{\raggedright\arraybackslash}p{(\linewidth - 10\tabcolsep) * \real{0.1364}}
  >{\raggedleft\arraybackslash}p{(\linewidth - 10\tabcolsep) * \real{0.1818}}
  >{\raggedleft\arraybackslash}p{(\linewidth - 10\tabcolsep) * \real{0.1818}}
  >{\raggedleft\arraybackslash}p{(\linewidth - 10\tabcolsep) * \real{0.1818}}
  >{\raggedleft\arraybackslash}p{(\linewidth - 10\tabcolsep) * \real{0.1818}}
  >{\centering\arraybackslash}p{(\linewidth - 10\tabcolsep) * \real{0.1364}}@{}}
\toprule\noalign{}
\begin{minipage}[b]{\linewidth}\raggedright
benchmark × model
\end{minipage} & \begin{minipage}[b]{\linewidth}\raggedleft
Mem0
\end{minipage} & \begin{minipage}[b]{\linewidth}\raggedleft
KDF
\end{minipage} & \begin{minipage}[b]{\linewidth}\raggedleft
gap
\end{minipage} & \begin{minipage}[b]{\linewidth}\raggedleft
\(p\)
\end{minipage} & \begin{minipage}[b]{\linewidth}\centering
winner
\end{minipage} \\
\midrule\noalign{}
\endhead
\bottomrule\noalign{}
\endlastfoot
LongMemEval 500 Q × gpt-4o-mini & 0.672 & 0.434 & \(-0.238\) & \(< 10^{-16}\) & Mem0 \\
LongMemEval 500 Q × gpt-4.1-mini & 0.722 & 0.452 & \(-0.270\) & \(3 \times 10^{-23}\) & Mem0 \\
LoCoMo temporal 321 Q × gpt-4o-mini & 0.206 & 0.312 & \(+0.106\) & \(1.4 \times 10^{-3}\) & KDF \\
LoCoMo temporal 321 Q × gpt-4.1-mini & 0.090 & 0.324 & \(+0.234\) & \(1.6 \times 10^{-14}\) & KDF \\
\end{longtable}
}

→ \textbf{The benchmark-dependent division of labor is robust across models.} LLM-based memory (Mem0) is strong on short-dialogue general QA but has a structural weakness for long-conversation date/time recall. KDF is positioned as a dedicated complementary layer for the latter.

\textbf{F-061--F-065 --- Domain-generalization experiments (2026-04-19, zero added cost)}:

F-044--F-060 validates only the LLM-memory domain. To position KDF as a \textbf{general-purpose graph compression technique}, we verified on four additional domains:

{\def\LTcaptype{none} % do not increment counter
\begin{longtable}[]{@{}
  >{\raggedright\arraybackslash}p{(\linewidth - 6\tabcolsep) * \real{0.2500}}
  >{\raggedright\arraybackslash}p{(\linewidth - 6\tabcolsep) * \real{0.2500}}
  >{\raggedright\arraybackslash}p{(\linewidth - 6\tabcolsep) * \real{0.2500}}
  >{\raggedright\arraybackslash}p{(\linewidth - 6\tabcolsep) * \real{0.2500}}@{}}
\toprule\noalign{}
\begin{minipage}[b]{\linewidth}\raggedright
Finding
\end{minipage} & \begin{minipage}[b]{\linewidth}\raggedright
Task
\end{minipage} & \begin{minipage}[b]{\linewidth}\raggedright
Result
\end{minipage} & \begin{minipage}[b]{\linewidth}\raggedright
Implication
\end{minipage} \\
\midrule\noalign{}
\endhead
\bottomrule\noalign{}
\endlastfoot
\textbf{F-061} & Betweenness / APSP on 4 synthetic graphs & \textbf{Mixed}: KDF wins on ER / SBM, loses BA / WS for betweenness; wins APSP on all 4 for path distance & TopDegree wins on scale-free graphs; KDF wins on uniform / community graphs \\
\textbf{F-062} & Git-commit pruning (tokio, 4752 commits) & \textbf{Positive}: merge recall \(99.5\%\) at \(30\%\) keep & Validated for commercial git archival \\
\textbf{F-063} & Gaussian Process inducing points (California housing, Friedman1) & \textbf{Negative}: KDF \textless{} Random \textless{} KMeans on GP fit & Density estimation is not applicable to KDF \\
\textbf{F-064} & Naive call-graph API preservation (flask, Python \texttt{ast}) & \textbf{Negative}: KDF \(16\%\) vs.~Random \(41\%\) & Public APIs have high in-degree --- opposite of KDF Rare \\
\textbf{F-065} & Git-pruning 3-repo replication (tokio / pytest / lodash) & \textbf{Partial}: merge recall depends on the repo's merge rate (tokio \(99\%\) / pytest \(59\%\) / lodash \(100\%\)) & KDF matches TopDegree only when ``merges are rare'' in the repo \\
\end{longtable}
}

Combining these negative findings, \textbf{we establish the decisive predictor of KDF applicability}: ``does structural rareness correlate with task importance?''. LLM-memory temporal recall, path-based algorithms, and OSS-library-style git repos satisfy the correlation → KDF is effective. Density estimation, API-boundary detection, and merge-heavy enterprise repos do not → KDF is not applicable.

\textbf{Phase X --- Claim-level realistic benchmarks (2026-04-19, zero added cost)}:

After F-068 brought all three mechanisms of Claim 1 to empirically-backed status, we ran realistic benchmarks for the remaining claim groups as Phase X:

{\def\LTcaptype{none} % do not increment counter
\begin{longtable}[]{@{}
  >{\raggedright\arraybackslash}p{(\linewidth - 6\tabcolsep) * \real{0.2500}}
  >{\raggedright\arraybackslash}p{(\linewidth - 6\tabcolsep) * \real{0.2500}}
  >{\raggedright\arraybackslash}p{(\linewidth - 6\tabcolsep) * \real{0.2500}}
  >{\raggedright\arraybackslash}p{(\linewidth - 6\tabcolsep) * \real{0.2500}}@{}}
\toprule\noalign{}
\begin{minipage}[b]{\linewidth}\raggedright
Finding
\end{minipage} & \begin{minipage}[b]{\linewidth}\raggedright
Task
\end{minipage} & \begin{minipage}[b]{\linewidth}\raggedright
Result
\end{minipage} & \begin{minipage}[b]{\linewidth}\raggedright
Implication
\end{minipage} \\
\midrule\noalign{}
\endhead
\bottomrule\noalign{}
\endlastfoot
\textbf{F-069} & Claim 5 / 14 / 17 on LoCoMo temporal 321 Q & \textbf{Mixed}: C17 distributed-execution bit-exact pass; C5 / C14 inferior to KDF\_static on static tasks (all time-aware conditions have \(\Delta \le 0\)) & Time signals are subsumed by structural rareness and redundant on static tasks; streaming benefit is empirical for \emph{temporally recurring rare} (F-072 NASA \(+3.06\) pt) but absent / inverted for \emph{one-shot rare} \textbf{{[}v2: F-087 Apache \(-13.04\) pt{]}} \\
\textbf{F-070} & Claim 36--41 T\_wait + Claim 47--48 sandwich, realistic & \textbf{Mixed}: mechanism {[}True{]} / canonical values {[}False{]}. Four-benchmark refutation of \((0.70, 0.80)\), F1 \(0.000\) vs.~F1\(((0.70, 1.00)) = 1.000\); LoCoMo streaming \(100\%\) RARE demoted under canonical & The 2-threshold \emph{mechanism} is maintained as novel; specific canonical values require domain-specific calibration \\
\textbf{F-071} & Claim 20--32 dynamic control on LoCoMo streaming (lightweight) & \textbf{Mechanism {[}True{]} / no selection benefit}: Claim 21 \(5:3:1\) integer tick exact; MetaController \(\alpha\)-bound clamp working; TransitionController ceiling-effected (F-031 confirmed) & Mechanism-only validation; the streaming value is anchored on NASA-type \emph{recurring rare} (F-072) but does not generalize to \emph{one-shot rare} \textbf{{[}v2: F-087{]}} \\
\textbf{F-072} & Claim 14 / 25 / 27--32 on NASA HTTP 50 k streaming replay & \textbf{{[}True{]} Claim 14 \(+3.06\) pt} over static baseline (C1 decay \(0.4898\) vs.~C0 \(0.4592\)); {[}Warning{]} Claim 25 activation neutralizes on evenly distributed rare; Claim 27--32 are selection-neutral (as predicted) & \textbf{The paper narrowing ``streaming is the true use case'' is empirically validated.} The first positive evidence for Claim 14's value proposition. ActivationScore requires discrimination of the rare-event temporal pattern. \textbf{{[}v2{]}} F-087 Apache replication (\(-13.04\) pt) further narrows this anchor to \emph{temporally recurring rare events}; for one-shot rare, streaming with decay is actively harmful (sign reversal) \\
\end{longtable}
}

These Phase X findings both strengthen paper credibility by ``self-refuting our own canonical values on four benchmarks'' and, via F-072, provide a \textbf{positive empirical anchor after the narrowing}. Two follow-up directions are explicitly indicated as future work: ``domain-calibrated parameter auto-tuning'', and ``conditional use of Claim 25 activation based on the rare-event temporal pattern''.

\textbf{Phase 2 / Phase 2.5 --- Scope refinement after v1 release (2026-04-20 -- 2026-04-29)}:

Following Zenodo v1 publication, we systematically tested the v1 claims
against new domains. The 23 findings (F-073 -- F-100) \textbf{sharpen} rather
than overturn the v1 narrative: they confirm that the ``decisive
predictor'' framework articulated in §6.1 (does structural rareness
correlate with task importance?) is the right level of abstraction, and
they delineate the resulting \textbf{7-pattern empirical arc} (5 narrowing
findings + 2 positive cross-domain replications).

{\def\LTcaptype{none} % do not increment counter
\begin{longtable}[]{@{}
  >{\raggedright\arraybackslash}p{(\linewidth - 6\tabcolsep) * \real{0.2500}}
  >{\raggedright\arraybackslash}p{(\linewidth - 6\tabcolsep) * \real{0.2500}}
  >{\raggedright\arraybackslash}p{(\linewidth - 6\tabcolsep) * \real{0.2500}}
  >{\raggedright\arraybackslash}p{(\linewidth - 6\tabcolsep) * \real{0.2500}}@{}}
\toprule\noalign{}
\begin{minipage}[b]{\linewidth}\raggedright
Finding
\end{minipage} & \begin{minipage}[b]{\linewidth}\raggedright
Task
\end{minipage} & \begin{minipage}[b]{\linewidth}\raggedright
Result
\end{minipage} & \begin{minipage}[b]{\linewidth}\raggedright
Implication
\end{minipage} \\
\midrule\noalign{}
\endhead
\bottomrule\noalign{}
\endlastfoot
\textbf{F-073} & Wikipedia orphan article preservation (simplewiki, \(\rho=0.20\)) & KDF \(-4.07\) pt vs.~Random & scale-free orphan pool with KDF Rare layer protecting deg=1 nodes is \textbf{wrong direction} for orphan importance; predictable scope violation \\
\textbf{F-074} & BGL supercomputer anomaly templates (\(\rho=0.20\)) & KDF \(-12.92\) pt & anomaly templates are \textbf{hub-like}, KDF Rare layer protects opposite; predictable scope violation \\
\textbf{F-075} & Citation interdisciplinary bridge detection (\(\rho=0.20\)) & KDF complete loss (recall \(0\%\)) & bridges occupy \textbf{Core-like} structural position; KDF Rare layer cannot capture; predictable scope violation \\
\textbf{F-076 / F-077} & Git archival cross-repo expansion (N=3 \(\to\) N=9 OSS repos) & merge-rate threshold pattern decisively confirmed; vscode-style high-merge-rate repos remain exceptions & F-062 / F-065 \textbf{strengthens} as a Tier-A product on OSS-style repos with merge rate \(< 10\%\) \\
\textbf{F-082} + \textbf{F-085 Task 1} & MovieLens niche genres (6 genres × bipartite, k=1 default) & γ-check \(100\%\) across all 6 genres (Film-Noir / IMAX / Western / Musical / War / Documentary) & Niche content surfacing \textbf{strengthens} as a Tier-A product, genre-agnostic robustness \\
\textbf{F-085 Task 4} & Reddit community anomaly title replication & replication failure & Reddit community anomaly product candidate \textbf{withdrawn} \\
\textbf{F-084} & Biological PPI + OncoKB cancer gene (\(N=1077\)) & γ-check \(74\% < 95\%\) Strong threshold; cancer genes are hub-biased (TP53/BRCA1/MYC) & Biological domain reject under hub-biased label; scope boundary refined \\
\textbf{F-086 α} & Real cgb-kdf on Git commit bipartite (flask, prettier) & merges land in \textbf{Core layer}, not Rare & Git archival product \textbf{reframed} as Core preservation, not Rare identification (framework correction) \\
\textbf{F-086 β} & ogbn\_arxiv academic citation meta-family & reject (peer-network structure, citing papers are themselves low-deg) & Academic citation domain reject; F-075 generalizes \\
\textbf{F-086 γ} & Domain-fit predictor sweep (N=5 meta-families) & 3 PASS / 2 REJECT; hub-peripheral / hub-biased / peer-network distinction validated & \textbf{F-086 γ replaces \texttt{bias-detector} as the working predictor} (γ-check correlation rate framework) \\
\textbf{F-087} & Apache error log streaming replication (LogHub Apache.log, freq \(\le 10\) rare) & streaming \(-13.04\) pt vs.~C0 static; full streaming reaches recall \(0.000\) & \textbf{Streaming benefit (F-072) narrows to temporally recurring rare}; one-shot rare actively harmed by decay + activation \\
\textbf{F-090} & \texttt{bias-detector} predictor on N=21 datasets, certain prediction \(= 11\) & accuracy \(5/11 = 45.5\% \ll 70\%\) pre-registered threshold & \texttt{bias-detector} commercial path \textbf{withdrawn}; v1 §6.1 reference is updated; F-086 γ remains the working predictor \\
\textbf{F-091} & Claim 10 (\(\alpha_{\text{core}}=2.0\), ``core of the invention'') cross-domain robustness; \(\alpha \in \{0.5, 1.0, 2.0, 3.0, 4.0\}\) sweep on NASA + Apache streaming + MovieLens null control & NASA: \(\alpha=2.0\) optimal (diff \(0.00\) pt vs best); Apache: \(\alpha=4.0\) optimal, \(\alpha=2.0\) diff \(-17.39\) pt; MovieLens null PASS & \textbf{H\_α PARTIAL (1/2)}: Claim 10 mechanism supported, but canonical \(\alpha=2.0\) is NASA-recurring-rare-specific; one-shot rare requires domain-specific \(\alpha\) calibration. Sister to F-070 (sandwich canonical refute) \\
\textbf{F-092} & Claim 31 Lyapunov stability under real-data adversarial burst perturbation (NASA HTTP w50 + 1000 rare-target events) & Boundedness PASS (\(\alpha_{\text{edge}} \in [1.0, 2.5]\) all 100 windows); Recovery PASS (diff \(0.0043\) at w55); Functional FAIL (recall\_perturbed \(/\) recall\_baseline \(= 0.000 / 0.4592\)) & \textbf{H\_L PARTIAL (2/3)}: Controller mechanism robust, but adversarial degree inflation demotes natural rare resources from Rare to Core layer. \textbf{4th self-refutation pattern} (F-070 / F-087 / F-091 sister); production deploy needs upstream rate-limit / provenance defense layer \\
\textbf{F-093} & NASA F-072 anchor robustness to rare code subset; 5 variants × 5 conditions = 25 runs at \(\alpha_{\text{core}}=2.0\) canonical & V1 (4xx+5xx) \(\Delta = +3.06\) pt; V2 (4xx only) and V4 (404 only) identical \(+3.06\) pt; V3 (5xx only) and V5 (500 only) trivial (n\_rare \(= 0\), NASA log has zero 5xx responses) & \textbf{H\_R PASS (3/5) + anchor sharpening category} (distinct from self-refutation): F-072 anchor is \emph{substantively} ``rare = 404 page-not-found pattern driven'' on this dataset; KDF captures rare resource pattern, not rare code set. F-072 anchor specificity now sharpened on 3 axes (domain via F-087, \(\alpha\) via F-091, rare type via F-093). Generalization to 5xx-containing logs is future work \\
\textbf{F-094} & Apache recurring-rare positive replication of F-072 streaming benefit; same Apache dataset (LogHub Apache.log, 31,062 records) as F-087 with rare def flipped to recurring (freq \(\ge 5\)); 2 rare defs × 5 conditions = 10 runs at \(\alpha_{\text{core}}=2.0\) canonical & Sanity (V\_one-shot, F-087 reproduce): \(\Delta = -13.04\) pt \(\pm 0.0\) pt (preprocessing / build env consistent); Main (V\_recurring): \(\Delta = +3.67\) pt \(> +1.0\) pt threshold & \textbf{H\_R+ PASS} + \textbf{positive direction completion of the self-refutation arc}: F-072 anchor's true axis (recurring rare structure, hypothesized via F-093 structural reading) is empirically replicated cross-domain (NASA \(+3.06\) + Apache recurring \(+3.67\)). 5 patterns now (4 narrow + 1 positive); scope of Claim 14 streaming benefit narrow but durable on recurring-rare axis \$\times \alpha=2.0 \times \$ NASA-Apache-homotype log dataset \\
\textbf{F-097} & BGL HW kernel log recurring-rare cross-domain N=3 expansion of F-094; LogHub BGL Blue Gene/L (300,000-line subsample, anomaly subgraph 79,641 records, 8 unique content templates); recurring def freq \(\ge 5\) at \(\alpha_{\text{core}}=2.0\) canonical & Main (V\_recurring): \(\Delta_{\text{recurring}} = +33.33\) pt \(\gg +1.0\) pt threshold (C0=0.0 → C2/C4=0.333, 1/3 rare templates retained); Sanity (V\_one-shot): \(\Delta = +0.00\) pt inconsistent with F-087 / F-094 negative direction & \textbf{H\_R+ PASS} + \textbf{2nd positive replication, cross-family}: durable recurring-rare axis now N=3 (NASA HTTP + Apache error + BGL HW kernel), expanding from web access log family to HW kernel log family. \textbf{Sanity inconsistent narrows one-shot disposal mechanism to web log family with rich resource alphabet} --- BGL alphabet of 8 anomaly templates is too small for streaming discrimination room. Caveat: \(n_{\text{rare}} = 3\), recall granularity \(\{0, 1/3, 2/3, 1\}\), direction unambiguous at low-N power \\
\textbf{F-099} & v1 / v2 / v3 router design-space characterization across 5 cells × 3 variants (4 paid LongMemEval+LoCoMo + 1 local LongMemEval, deterministic post-hoc replay); v1 = precision-only, v2 = precision + length≥100, v3 = length-only & H\_v1\_paid PASS\_negative (LongMemEval paid both models: F-053 \(\Delta_{v1} = -11.60\) pt \(p=10^{-7}\), F-059 \(\Delta_{v1} = -13.00\) pt \(p=10^{-10}\)); H\_v1\_local REPRODUCED (F-096 LongMemEval local \(\Delta_{v1} = +6.26\) pt deterministic); Sanity 4/4 paid v2 cells match F-060 published values within \(\pm 0.0004\); v3 (length-only) matches v2 (\(+9.66 / +22.43\) pt) on LoCoMo cells & \textbf{F-099 anchors v2's length≥100 component necessity} via indirect logic (v1 -11/-13 pt vs v2 0 pt on paid LongMemEval): without the length filter, routing precision queries to KDF in paid setting is actively harmful. \textbf{Important caveat: F-099 does NOT establish precision component's independent necessity} --- v3 length-only matches v2 on LoCoMo, suggesting precision filter may be redundant once length filter holds. v1 is \textbf{not a product candidate}. First test case of pre-reg template \texttt{\_template\_pre\_reg.md} \\
\textbf{F-100} & HDFS distributed file system log empirical verification of recurring-rare cross-domain N=4 expansion; Loghub HDFS\_v1 (100,000 BlockId subsample, 2.35M edges, 26 templates, 4.90\% anomaly); 3 hypotheses (H\_R+\_literal recurring, H\_anomaly HDFS-native, Sanity V\_one-shot) at \(\alpha_{\text{core}}=2.0\) canonical & H\_R+\_literal \(\Delta = -23.08\) pt \textbf{FAIL} (recurring rare benefit refuted on HDFS); H\_anomaly \(\Delta = -4.30\) pt \textbf{FAIL} (HDFS-native anomaly preservation FAIL); Sanity \(\Delta = +0\) pt inconsistent (one-shot disposal absent, reinforcing F-097 small-alphabet caveat) & \textbf{Both empirically FAILED} + \textbf{5th narrowing pattern}: cross-domain durability arc narrowed to \textbf{web access log family + HW kernel log family} specific (NASA HTTP + Apache error log + BGL HW kernel) --- distributed file system log domain NOT covered by current claims. N=4 not achieved. 6-pattern arc → \textbf{7-pattern arc} (5 narrow + 2 positive). Pre-reg template second high-severity test case: §0.1 anchor constraint predicted small-alphabet caveat would re-confirm → empirically borne out; §0.3 observation/interpretation separation prevented narrative drift under both-FAIL outcome \\
\end{longtable}
}

The combined picture: KDF's empirically validated core is now \textbf{four
product candidates} (Obsidian / MovieLens niche / Mem0 hybrid Router /
Git Core preservation) plus the F-086 γ domain-fit framework. The v1
``narrow but durable'' theme is preserved; v2 sharpens the empirical
boundary at the cost of three specific positionings (direct
SOTA-beating, \texttt{bias-detector} as commercial predictor, and Claim 10 /
Claim 14 / Claim 31 universal value claims --- all narrowed to
domain-conditional or scope-conditional form via the \textbf{five-finding
self-refutation pattern} F-070 / F-087 / F-091 / F-092 / F-100). \textbf{F-093
(v2.1) adds an ``anchor sharpening'' category distinct from
self-refutation}: the F-072 anchor is now resolved at higher
resolution as ``NASA-recurring-rare \$\times \alpha=2.0 \times \$
404-pattern driven''. \textbf{F-094 (v2.2) opens the positive direction of
the arc}: the recurring-rare axis hypothesized via F-093 structural
reading is empirically replicated cross-domain (\(+3.67\) pt on Apache
same dataset with rare def flipped to freq \(\ge 5\), sanity F-087
reproduce \(\pm 0.0\) pt). \textbf{F-097 (v2.5) extends positive replication to
N=3 cross-family} (BGL HW kernel log \(+33.33\) pt), and \textbf{F-100 (v2.7)
attempts but fails the N=4 expansion} to distributed file system log
(HDFS H\_R+\_literal \(-23.08\) pt FAIL, H\_anomaly \(-4.30\) pt FAIL),
narrowing cross-family durability to web access log family + HW kernel
log family specific. The full arc is now \textbf{7 patterns total = 5
narrowing + 2 positive}; Claim 14 streaming benefit is empirically
supported on N=3 datasets (NASA HTTP + Apache error + BGL HW kernel)
when restricted to recurring rare in those two log families.
Generalization to distributed file system log family is empirically
ruled out for current claims, and other log families remain future
work. See \href{../PHASE_2_RETROSPECTIVE.md}{\texttt{docs/PHASE\_2\_RETROSPECTIVE.md}}
for the single-document retrospective.

\textbf{Foreign baseline anchor (Mem0 family, N=3 empirical cells)}:

KDF vs.~the Mem0 family of LLM-memory systems is empirically anchored
across three operational settings, all completed prior to the F-094
arc completion:

\begin{enumerate}
\def\labelenumi{\arabic{enumi}.}
\tightlist
\item
  \textbf{F-048 (2026-04-18, local LLM proxy, zero cost)}: Local
  Mem0-style pipeline (Qwen2.5-0.5B fact extraction + BGE-small
  embedding) on LongMemEval first-100 questions yielded recall
  \(0.5083\) vs.~KDF \(0.8210\) (\textbf{KDF \(+31.27\) pt}). Under local /
  budget / privacy constrained environments, KDF retrieval is
  decisively superior to local Mem0-style approaches.
\item
  \textbf{F-053 (2026-04-18, paid Mem0 framework with gpt-4o-mini)}:
  Real KDF vs.~paid Mem0 on LongMemEval 500 Q --- Mem0 \(0.672\) vs.
  KDF \(0.434\) (\textbf{Mem0 \(+23.8\) pt}, \(p < 10^{-16}\), Mem0 wins
  significantly in 5 of 6 categories). On standard paid LLM-memory
  workloads, Mem0 is the stronger standalone system.
\item
  \textbf{F-060 (2026-04-19, KDF + Mem0 Router, zero added cost, detailed
  below)}: KDF + Mem0 hybrid Router (precision-query regex +
  conversation length \(\ge 100\) turns) is \textbf{strictly better than Mem0
  alone across 4 cells × 2 models} (\(+10\) to \(+23\) pt on LoCoMo
  temporal; \(\pm 0\) pt on LongMemEval; never worse).
\end{enumerate}

These three cells establish the Mem0-family foreign baseline at \(N=3\).
The position is \textbf{setting-dependent}: KDF favored in local / budget
environments (F-048), Mem0 favored as standalone in standard paid
LLM-memory workloads (F-053), and the \textbf{KDF + Mem0 hybrid Router} is
strictly better than Mem0 alone in hybrid deployments (F-060) ---
placing KDF as a \textbf{Mem0 complementary layer}, not a replacement, with
the Router implementation as the validated commercial path. The
remaining Cat-5 baselines (Letta / Zep / GraphRAG / Anthropic memory /
OpenAI memory / LangChain) are deferred to a paid-budget multi-session
sprint.

\textbf{Attempted local replication of F-060 (F-095 / F-096, 2026-04-29, infrastructure honest record)}:
A 4th cell aimed to replicate F-060's hybrid Router benefit in a fully
local environment (Ollama-served LLM + HuggingFace BGE embedder + Qdrant
local store, zero API cost). Two attempts:
F-095 (\texttt{llama3.1:8b-instruct-q4\_K\_M}, default ingest batching) was stopped
at the 5-question checkpoint due to wall-clock infeasibility (extrapolated
\(\sim 33\) -- \(36\) hours for the full LongMemEval + LoCoMo suite); the
root cause was Mem0 framework's two-LLM-call-per-\texttt{add()} cost
multiplied by small ingestion batches. F-096 (\texttt{qwen2.5:3b-instruct-q4\_K\_M}
with single-add ingest for LongMemEval and batch=50 for LoCoMo) reduced
wall-clock to \(\sim 2.5\) hours and completed both benchmarks (LongMemEval
\(n=479\), LoCoMo \(n=321\)), but the result is \textbf{methodologically inconclusive}:
LoCoMo Mem0 alone scored \(0/321\) correct and KDF alone scored \(2/321\)
under the local 3B answer-generator and judge --- a sub-noise floor that
prevents meaningful \(\Delta_{\text{router}}\) measurement (mechanical
verdict: H\_R PARTIAL with \(\Delta = +0.62\) pt, \(p = 0.5\)). Sanity 2
local-vs-paid baseline shifts of \(-20\) to \(-64\) pt across all four cells
indicate that qwen2.5-3b is below the threshold to substitute for paid
gpt-4o-mini on these benchmarks. F-095 / F-096 do \textbf{not refute F-060} ---
they record an infrastructure threshold instead. A valid local replication
will require a stronger local LLM (\(7\)B+ class) and is left as
future work. A descriptive cross-LLM-size finding does emerge: the Mem0
framework's batched fact compression destroys \(\sim 99\%\) of raw turn
substrings (\texttt{mem0\_recall\_substring} \(= 0.000\) -- \(0.006\) across F-095
8B and F-096 3B), in contrast to F-048's hand-rolled per-turn extraction
(recall \(= 0.508\)); this sharpens F-048's ``weak LLM degrades retrieval''
caveat into ``framework compression strategy itself does not preserve
substrings, independent of LLM size''.

\textbf{F-060 --- Empirical validation of a complementary architecture via the Ext-1 Precision-Query Router (2026-04-19, zero added cost)}:

Applying the following routing logic post-hoc to the existing data of F-053 / 057 / 058 / 059:

\begin{verbatim}
if is_precision_query(q) and conversation_length >= 100 turns:
    use KDF answer
else:
    use Mem0 answer
\end{verbatim}

Result (v2 = precision + long context):

{\def\LTcaptype{none} % do not increment counter
\begin{longtable}[]{@{}
  >{\raggedright\arraybackslash}p{(\linewidth - 8\tabcolsep) * \real{0.1579}}
  >{\raggedleft\arraybackslash}p{(\linewidth - 8\tabcolsep) * \real{0.2105}}
  >{\raggedleft\arraybackslash}p{(\linewidth - 8\tabcolsep) * \real{0.2105}}
  >{\raggedleft\arraybackslash}p{(\linewidth - 8\tabcolsep) * \real{0.2105}}
  >{\raggedleft\arraybackslash}p{(\linewidth - 8\tabcolsep) * \real{0.2105}}@{}}
\toprule\noalign{}
\begin{minipage}[b]{\linewidth}\raggedright
cell
\end{minipage} & \begin{minipage}[b]{\linewidth}\raggedleft
Mem0 alone
\end{minipage} & \begin{minipage}[b]{\linewidth}\raggedleft
Router (v2)
\end{minipage} & \begin{minipage}[b]{\linewidth}\raggedleft
gain
\end{minipage} & \begin{minipage}[b]{\linewidth}\raggedleft
\(p\)
\end{minipage} \\
\midrule\noalign{}
\endhead
\bottomrule\noalign{}
\endlastfoot
LongMemEval 500 Q × gpt-4o-mini & 0.672 & 0.672 & 0.000 & 1.00 (safe) \\
LongMemEval 500 Q × gpt-4.1-mini & 0.722 & 0.722 & 0.000 & 1.00 (safe) \\
LoCoMo temporal 321 Q × gpt-4o-mini & 0.206 & \textbf{0.302} & \textbf{\(+0.097\)} & 0.003 {[}Significant{]} \\
LoCoMo temporal 321 Q × gpt-4.1-mini & 0.090 & \textbf{0.315} & \textbf{\(+0.224\)} & \(4 \times 10^{-14}\) {[}Significant{]} \\
\end{longtable}
}

\emph{Note on precision: the columns ``Mem0 alone'' and ``Router (v2)'' are 3-decimal rounded for readability; the ``gain'' column is computed from the 4-decimal raw scores and may therefore differ from naive subtraction of the displayed values by ±0.001.}

The Router is \textbf{greater than or equal to Mem0 alone across all four cells} (strictly-better property); on long-conversation precision queries, the improvement reaches \(+22.4\) pt.~LLM-API call count is reduced by \(97\%\) on long conversations. This is the first quantitative validation of this paper's architectural thesis that KDF should be designed as \textbf{a complementary layer to Mem0, not a replacement}.

\textbf{{[}v2.6{]}} F-099 router design-space characterization (5 cells × 3 variants, deterministic post-hoc replay) anchors the \textbf{necessity of v2's length≥100 component}: the v1 variant (precision-only, no length filter) is significantly worse on paid LongMemEval (\(-11.6\) / \(-13.0\) pt, \(p < 10^{-7}\)) because routing precision queries to KDF in the paid setting (where Mem0 \(\gg\) KDF) is actively harmful without the length gate. The v3 variant (length-only, no precision filter) matches v2 (\(+9.66 / +22.43\) pt) on LoCoMo cells, suggesting the precision filter may be redundant once the length filter holds on these cells; precision filter independent necessity remains evaluated by other anchors, not by F-099. v1 is \textbf{not a product candidate} (paid hurt significant, local help at sub-noise floor in F-096).

\begin{itemize}
\tightlist
\item
  \textbf{KDF's benchmark-dependent behavior}: in long conversations (LoCoMo's 300--700 turns/conv, with date/time information scattered across raw turns), KDF's raw-turn retention is advantageous; in short dialogues (LongMemEval's 20--30 turns/Q), Mem0's fact extraction dominates.
\end{itemize}

Therefore, \textbf{KDF's realistic positioning} pivots to the following:

\begin{enumerate}
\def\labelenumi{\arabic{enumi}.}
\tightlist
\item
  \textbf{Long-conversation temporal-recall use cases} (F-056: \(+22\) pt validated) --- meeting minutes, month-scale journals, year-scale conversation histories where date/time information must be referenced; KDF's raw-turn retention is decisively advantageous.
\item
  \textbf{Environments where cost / latency / privacy / determinism take priority over accuracy} --- local-first chatbot, air-gapped agent, \(< 1\)-ms real-time memory gating, budget-constrained deployment, deterministic regulated output.
\item
  \textbf{Retention strategies smarter than TTL} (real 500 Q: KDF recall \(0.665\) vs.~TTL\_recent \(0.180\), \(\times 3.7\)).
\end{enumerate}

\textbf{Honest limits}: in the general LLM-memory market of LongMemEval type (short dialogue + concise questions), KDF does not reach Mem0's accuracy. KDF use in this market is limited to cases that tolerate the accuracy trade-off.

For other domains, we honestly declare what KDF can and cannot do, and provide it in an openly verifiable form. In particular, \textbf{KDF use in the general-retrieval market is explicitly discouraged} (F-045).

\begin{center}\rule{0.5\linewidth}{0.5pt}\end{center}

\section{8. Theoretical Foundation --- KDF as a Computational Realization of Structural Holes Theory}\label{theoretical-foundation-kdf-as-a-computational-realization-of-structural-holes-theory}

Having accumulated the empirical characteristics of KDF (F-061--F-065), we find that KDF's behavior \textbf{algorithmically aligns with the classical ``Structural Holes'' theory of organizational sociology \citep{burt1992structuralholes}}. This is not a post-hoc analogy; it refers to a \textbf{graph-theoretic isomorphism} (with explicit caveats developed in the \emph{Limitations and Risks of the Theoretical Claim} subsection at the end of §8; the isomorphism is a strong framing hypothesis, not a proved theorem).

\subsection{\texorpdfstring{Summary of Burt's Structural Holes theory \citep{burt1992structuralholes}}{Summary of Burt's Structural Holes theory {[}@burt1992structuralholes{]}}}\label{summary-of-burts-structural-holes-theory-burt1992structuralholes}

From an empirical study of workplace networks (\(n = 500\) MBA students), Burt argued:

\begin{enumerate}
\def\labelenumi{\arabic{enumi}.}
\tightlist
\item
  \textbf{The individuals who hold the highest information, innovation, and negotiation power in an organization are not the central figures within dense clusters, but the ``brokers'' who bridge different clusters.}
\item
  A \textbf{structural hole} is the empty space where two clusters are not directly connected.
\item
  By spanning this hole, the broker can:

  \begin{itemize}
  \tightlist
  \item
    control information non-redundancy (different information flows on each side),
  \item
    set exchange terms favorably (brokerage power),
  \item
    cross-pollinate ideas (innovation advantage).
  \end{itemize}
\item
  Brokers typically have \textbf{low degree} (fewer ties make bridging more efficient) and \textbf{high betweenness} (they lie on paths).
\end{enumerate}

\subsection{Mathematical Alignment with KDF's Behavior}\label{mathematical-alignment-with-kdfs-behavior}

KDF's Rare / Core / Edge / Garbage classification corresponds to Burt's broker concept as follows:

{\def\LTcaptype{none} % do not increment counter
\begin{longtable}[]{@{}
  >{\raggedright\arraybackslash}p{(\linewidth - 6\tabcolsep) * \real{0.2500}}
  >{\raggedright\arraybackslash}p{(\linewidth - 6\tabcolsep) * \real{0.2500}}
  >{\raggedright\arraybackslash}p{(\linewidth - 6\tabcolsep) * \real{0.2500}}
  >{\raggedright\arraybackslash}p{(\linewidth - 6\tabcolsep) * \real{0.2500}}@{}}
\toprule\noalign{}
\begin{minipage}[b]{\linewidth}\raggedright
KDF layer
\end{minipage} & \begin{minipage}[b]{\linewidth}\raggedright
Structural feature
\end{minipage} & \begin{minipage}[b]{\linewidth}\raggedright
Burt's classification
\end{minipage} & \begin{minipage}[b]{\linewidth}\raggedright
Evidence
\end{minipage} \\
\midrule\noalign{}
\endhead
\bottomrule\noalign{}
\endlastfoot
\textbf{Rare} (deg \(= 1\)) & Boundary node, unique connection & Pure broker (if between clusters) & F-012 Obsidian orphan \\
\textbf{Core} & Bridges many clusters, moderate degree & Structural broker & F-062 merge commits (\(99.5\%\) recall) \\
Edge & Within-cluster center, high degree & Cluster insider (non-broker) & F-061 scale-free hubs (KDF loses) \\
Garbage & Redundant, out of capacity & Peripheral / replaceable & F-052 low-recall answer turns \\
\end{longtable}
}

KDF's selection principle of \textbf{``prioritizing Rare + Core for protection''} is equivalent to \textbf{``prioritizing nodes occupying broker positions for protection''}. This is not a coincidence: the graph-theoretic definition of structural rareness (low degree and/or high betweenness potential) overlaps with Burt's broker definition.

\subsection{Empirical Validation via the Scale-Free vs.~Community-Graph Split}\label{empirical-validation-via-the-scale-free-vs.-community-graph-split}

The betweenness-centrality results of F-061 on four graphs (ER, BA, WS, SBM) constitute a \textbf{predictive validation} of this theoretical correspondence:

{\def\LTcaptype{none} % do not increment counter
\begin{longtable}[]{@{}
  >{\raggedright\arraybackslash}p{(\linewidth - 6\tabcolsep) * \real{0.2500}}
  >{\raggedright\arraybackslash}p{(\linewidth - 6\tabcolsep) * \real{0.2500}}
  >{\raggedright\arraybackslash}p{(\linewidth - 6\tabcolsep) * \real{0.2500}}
  >{\raggedright\arraybackslash}p{(\linewidth - 6\tabcolsep) * \real{0.2500}}@{}}
\toprule\noalign{}
\begin{minipage}[b]{\linewidth}\raggedright
Graph type
\end{minipage} & \begin{minipage}[b]{\linewidth}\raggedright
Structural feature
\end{minipage} & \begin{minipage}[b]{\linewidth}\raggedright
Burt-theoretic prediction
\end{minipage} & \begin{minipage}[b]{\linewidth}\raggedright
KDF result
\end{minipage} \\
\midrule\noalign{}
\endhead
\bottomrule\noalign{}
\endlastfoot
ER (Erdős--Rényi) & Uniform random & Brokers exist; bridge protection wins & {[}True{]} KDF wins (top-50 recall \(0.70\)) \\
SBM (Stochastic Block Model) & Planted communities & Inter-community brokers are decisive; KDF applicability is maximal & {[}True{]} KDF wins (\(0.50\) vs.~\(0.36\) TopDegree) \\
BA (Barabási--Albert, scale-free) & Hub-dominated & Hubs have abundant alternative paths → brokerage power is distributed; KDF at disadvantage & {[}False{]} KDF loses to TopDegree \\
WS (Watts--Strogatz, small world) & Clustered + shortcut & Shortcuts are broker-like but of moderate degree; complex & {[}False{]} KDF loses to TopDegree on betweenness \\
\end{longtable}
}

\textbf{KDF's loss on scale-free graphs} is consistent with Burt's theoretical prediction that \textbf{``the relative value of brokers declines in hub-dominated networks.''} Hubs have many alternative paths, so removing any given hub leaves the network robust. Brokerage power concentrates in inter-community bridges, but is distributed in pure scale-free networks.

\subsection{Connection to Game-Theoretic Brokerage Power}\label{connection-to-game-theoretic-brokerage-power}

In network-bargaining theory \citep{myerson1977graphs, calvoarmengol2004networks}, when two players A and B wish to transact but are not directly connected and must go through an intermediary C, \textbf{C's payoff is inversely proportional to the number of alternative paths}. This is the mathematical formulation of structural holes.

KDF can be interpreted as an algorithm that deterministically extracts the \textbf{``intermediaries with no alternative paths'' (monopolistic brokers)}. In F-061's APSP experiment, KDF's \(4 \times\) improvement over Random in coverage maintenance --- particularly on \textbf{Watts--Strogatz} (small world, shortcut-broker structure) --- is an instance of this theoretical prediction.

\subsection{Computational Complexity of the Implementation Algorithm}\label{computational-complexity-of-the-implementation-algorithm}

Burt's structural-holes computation is traditionally performed via \textbf{effective-size} / \textbf{constraint} metrics (\(O(V^2)\) or worse). KDF's Rare / Core / Edge / Garbage classification achieves approximately the same broker detection in \textbf{\(O(V + E)\)} structural counting + classification (exact equivalence is left for future work).

\textbf{Implication}: KDF may be the first graph algorithm that \textbf{detects structural holes / brokerage power deterministically in near-linear time}. In the 30+ years since Burt's 1992 work \citep{burt1992structuralholes}, structural-holes theory has been continuously cited in social-science research (M\&A analysis, supply-chain resilience, innovation diffusion), but \textbf{no scalable computational algorithm has been established}. KDF's \(O(V + E)\) classifier may remove the rate-limiting barrier to practical deployment.

\subsection{Applications Space (with Evidence Sketches)}\label{applications-space-with-evidence-sketches}

From this theoretical correspondence, KDF's commercial and research application space is as follows:

\begin{enumerate}
\def\labelenumi{\arabic{enumi}.}
\tightlist
\item
  \textbf{Enterprise network analytics}

  \begin{itemize}
  \tightlist
  \item
    Internal communication analysis (Slack / Teams) → identifying cross-team brokers
  \item
    M\&A target selection: extracting the ``monopolistic intermediary companies'' that should be acquired
  \item
    Supply chain: identifying BCP-critical, low-degree Tier 3/4 suppliers
  \end{itemize}
\item
  \textbf{Urban / logistics networks}

  \begin{itemize}
  \tightlist
  \item
    Disaster-time APSP preprocessing (validated in F-061)
  \item
    Critical-intersection identification in transportation networks
  \end{itemize}
\item
  \textbf{Cybersecurity}

  \begin{itemize}
  \tightlist
  \item
    Lateral-movement detection via ``rare inter-segment connections'' (additional measurement needed)
  \end{itemize}
\item
  \textbf{Enterprise archival}

  \begin{itemize}
  \tightlist
  \item
    Decision preservation via cross-thread bridge capture (homologous to F-062 git merge)
  \end{itemize}
\end{enumerate}

\subsection{Limitations and Risks of the Theoretical Claim}\label{limitations-and-risks-of-the-theoretical-claim}

For honesty, we make the following explicit:

\begin{itemize}
\tightlist
\item
  \textbf{Burt's theory is primarily validated on human organizational networks (\(< 10^3\) nodes)}; application at millions of nodes remains an empirical question.
\item
  KDF's Rare / Core classification is \textbf{closer to a necessary-condition approximation than a sufficient condition} for broker detection (as suggested by KDF's loss on BA graphs in F-061).
\item
  Structural-holes theory itself has a known ``weak on hub-dominated networks'' limitation; see critiques such as \citet{powell2005network}.
\item
  \textbf{Whether KDF's \(O(V + E)\) complexity is formally equivalent to Burt's effective-size metric is unproven}; we have only approximation and empirical agreement.
\end{itemize}

With these caveats maintained, we position \textbf{KDF as a strong candidate for the computational realization of structural-holes detection}.

\subsection{8.6 Position vs Modern Alternatives}\label{position-vs-modern-alternatives}

Two recent currents threaten to absorb the conceptual space KDF occupies. We identify each and clarify why KDF is not subsumed.

\textbf{vs Graph-Attention-Network brokerage detection.} A 2024--2025 line of work applies GAT-style architectures to community detection and brokerage identification --- for example GATFELPA \citep{gatfelpa2024}, unsupervised GAT autoencoders \citep{gat-autoencoder2025}, and three-view multi-attention community detection \citep{tkdd-3672081}. These approaches achieve strong modularity scores, but in the process they introduce three properties KDF deliberately avoids: \textbf{non-determinism} (attention weights are sampled / SGD-trained, not analytically derived), \textbf{retraining dependence} (parameters must be re-fitted for each new graph), and \textbf{unauditability} (the Core/Rare boundary is implicit in learned weights rather than a computed predicate). Burt's \citep{burt1992holes} original constraint measure was prized precisely because it was analytically tractable and therefore reportable in a sociology paper without re-execution. KDF preserves that property --- its layer assignments are reproducible bit-exactly across runs (Claim 15) and can be re-derived by an auditor without access to training data --- while extending the framework with temporal decay (Claim 14) that classical structural-holes theory does not address. Thus KDF should be read as a \textbf{deterministic, audit-grade complement} to GAT methods, not a competitor on raw modularity.

\textbf{vs long-context LLMs that obviate memory.} Anthropic Claude with 1M-token context and Gemini 3 Pro with 10M tokens have prompted claims that explicit memory layers are obsolete. Empirically this is not borne out: the Chroma ``Context Rot'' study \citep{chromacontextrot2025} documents systematic reasoning degradation past roughly 100K tokens regardless of model class, and Anthropic's own context-length analysis confirms the trend. KDF's role in this landscape is specific: \textbf{deterministic selection of N items into a window where the LLM still reasons reliably}. Embedding-based retrieval stacks (Mem0, Letta, Zep) cannot offer this guarantee --- their selection is non-deterministic and varies with the underlying model. KDF's value is therefore not as a memory replacement (a thesis the F-053 LongMemEval comparison disposed of) but as an \textbf{audit-grade pre-filter} that operates upstream of retrieval. The F-060 Ext-1 router result (Mem0 + KDF strictly \textgreater{} Mem0 alone on date/time literal recall) is the empirical anchor for this position.

\begin{center}\rule{0.5\linewidth}{0.5pt}\end{center}

\section*{Acknowledgments}\label{acknowledgments}

We ran independent verification agents (GPT- and Claude-based) at 12 phase boundaries to check the validity of both the positive and negative claims of this work. AI collaboration was used at every stage --- specification freezing, code implementation, real-data verification, related-work survey, and drafting of this paper --- and the full verification process is recorded in \href{../VERIFIED_FINDINGS.md}{\texttt{docs/VERIFIED\_FINDINGS.md}}.

\begin{center}\rule{0.5\linewidth}{0.5pt}\end{center}

\section*{Appendix A: Summary of Key KDF Formulas}\label{appendix-a-summary-of-key-kdf-formulas}

{\def\LTcaptype{none} % do not increment counter
\begin{longtable}[]{@{}
  >{\raggedright\arraybackslash}p{(\linewidth - 4\tabcolsep) * \real{0.3333}}
  >{\raggedright\arraybackslash}p{(\linewidth - 4\tabcolsep) * \real{0.3333}}
  >{\raggedright\arraybackslash}p{(\linewidth - 4\tabcolsep) * \real{0.3333}}@{}}
\toprule\noalign{}
\begin{minipage}[b]{\linewidth}\raggedright
Claim
\end{minipage} & \begin{minipage}[b]{\linewidth}\raggedright
Formula
\end{minipage} & \begin{minipage}[b]{\linewidth}\raggedright
Description
\end{minipage} \\
\midrule\noalign{}
\endhead
\bottomrule\noalign{}
\endlastfoot
7 & \(C_{uv} = \deg(u) + \deg(v)\) & Local congestion \\
8, 9 & \(\lambda(C) = \beta(1 + \gamma C^\alpha)\), \(\alpha\) positive exponent & Nonlinear form of the decay rate (monotone increasing + power-law term) \\
10 & \(\alpha = 2\) & Fix the power-law exponent at 2 (core of the invention) \\
14 & \(w(t + dt) = w(t) \cdot \exp(-\lambda(C) \cdot dt)\) & Exponential decay law \\
15 & \(\deg_E(v) \le 1 \Rightarrow \text{Rare} \land \text{protected}\) & Rarity detection via absolute threshold \\
21 & \(dt_1 : dt_2 : dt_3 = 5 : 3 : 1\) & Update-period ratio across hierarchical regions \\
29 & \(\Delta \alpha \propto \delta k^4\) & Meta-control quartic law \\
44 & \(S_{\text{outer}} = \tfrac{7}{10} S_{\text{sys}} + \tfrac{2}{10} S_{\text{rel}} + \tfrac{1}{10} S_{\text{attr}}\) & Aggregated integrity score (outer; takes systematic / relational / attribute similarities as inputs) \\
45 & \(S_{\text{inner}} = 0.40 \cdot S_{\text{cos}} + 0.35 \cdot S_{\text{struct}} + 0.25 \cdot S_{\text{sign}}\) & Fingerprint similarity combination (inner; feeds into \(S_{\text{sys}}\) / \(S_{\text{rel}}\) / \(S_{\text{attr}}\) above) \\
46 & \(\phi(v) \in \mathbb{R}^{32}\), derived from Laplacian eigenvalues & Fixed-length structural fingerprint \\
47--48 & \(\theta_L = 0.70 \le S_{\text{outer}} \le \theta_U = 0.80\) & Sandwich acceptance condition, applied to the outer (aggregated) score from Claim 44 (canonical values --- refuted in §4.2; see §1.3 C3) \\
\end{longtable}
}

\begin{center}\rule{0.5\linewidth}{0.5pt}\end{center}

\section*{Appendix B: Implementation Architecture}\label{appendix-b-implementation-architecture}

\begin{itemize}
\tightlist
\item
  \texttt{crates/cgb-kdf/} --- reference implementation (Rust; PolyForm Noncommercial 1.0.0, commercial license separate); direct tests for all 50 patent claims.
\item
  \texttt{crates/bias-detector/} --- independently released crate; zero-dependency; a priori KDF applicability screening.
\item
  \texttt{demos/D1–D8/} --- showcase implementations across 8 domains.
\item
  \texttt{benchmarks/sota\_comparison/} --- SOTA comparison benchmarks.
\end{itemize}

GitHub: \href{https://github.com/ChaiCroquis/kdf-perovskite}{ChaiCroquis/kdf-perovskite}

\begin{center}\rule{0.5\linewidth}{0.5pt}\end{center}

\emph{Draft v0.3 --- 2026-04-19 (reflecting Phase X Step 1--5 completion). Comments and corrections welcome via GitHub issues.}

\begin{center}\rule{0.5\linewidth}{0.5pt}\end{center}


% -----------------------------------------------------------------------------
% Bibliography
% -----------------------------------------------------------------------------
\bibliography{references}

\end{document}
